\documentclass[11pt]{article}
\usepackage{series}

\renewcommand\SeriesName{Modal Triplet Theory}
\renewcommand\SeriesShort{Projection Duality}
\renewcommand\PartRoman{I}
\renewcommand\PartArabic{1}
\renewcommand\PartsInSeries{I}

\title{A Conditional Operator Encoding of Wave--Particle Duality in MTT\\
\large Local Kernels, Spectral Coherence, Instruments, and the Born-Bridge Boundary}
\author{Peter Nero}
\date{July 2026, Version 5}

% Paper-specific notation
\newcommand{\Kadm}{K_{\mathrm{adm}}}
\newcommand{\Badm}{B_{\mathrm{adm}}}
\newcommand{\ta}{\tau_{\mathrm{adm}}}
\newcommand{\epsadm}{\epsilon_{\mathrm{adm}}}
\newcommand{\Leff}{\Lambda_{\mathrm{eff}}}
\newcommand{\Hcoh}{\mathcal H_{\mathrm{coh}}}
\newcommand{\Hinc}{\mathcal H_{\mathrm{inc}}}
\newcommand{\Dab}{D_{ab}}
\newcommand{\Dmat}{D}
\newcommand{\Kdet}{K_{\mathrm{det}}}
\newcommand{\Efx}{E_x}
\newcommand{\calE}{\mathcal E}
\newcommand{\calD}{\mathcal D}
\newcommand{\calA}{\mathcal A}
\newcommand{\calR}{\mathcal R}
\newcommand{\calK}{\mathcal K}
\newcommand{\coh}{\mathrm{coh}}
\newcommand{\inc}{\mathrm{inc}}
\newcommand{\loc}{\mathrm{loc}}
\newcommand{\spec}{\mathrm{spec}}
\newcommand{\MTT}{\mathrm{MTT}}
\newcommand{\diag}{\mathrm{diag}}
\newcommand{\vis}{\mathrm{vis}}
\newcommand{\Meas}{\mathsf M}
\newcommand{\Basin}{\mathcal B}
\newcommand{\CND}{\mathrm{CND}}

% --- cleveref names for theorem-like environments ---
\crefname{assumption}{Assumption}{Assumptions}
\Crefname{assumption}{Assumption}{Assumptions}

\crefname{definition}{Definition}{Definitions}
\Crefname{definition}{Definition}{Definitions}

\crefname{proposition}{Proposition}{Propositions}
\Crefname{proposition}{Proposition}{Propositions}

\crefname{theorem}{Theorem}{Theorems}
\Crefname{theorem}{Theorem}{Theorems}

\crefname{corollary}{Corollary}{Corollaries}
\Crefname{corollary}{Corollary}{Corollaries}

\crefname{remark}{Remark}{Remarks}
\Crefname{remark}{Remark}{Remarks}

\begin{document}
\maketitle

\begin{abstract}
Wave--particle duality is usually presented as a primitive feature of quantum theory:
microscopic systems propagate as extended wave amplitudes, interfere through coherent
phase relations, and yet appear as localized particle-like events under measurement.
Modal Triplet Theory (MTT) proposes to encode this duality as a projection phenomenon:
one finite coherent-sector excitation has a local delta-like representation and a spectral
phase-coherent representation.

In a declared fixed-point regime, the relevant operator is
\[
  \Badm
  =
  P\chi(A)e^{-\ta A}\chi(A)P,
\]
where \(A\) is a stabilization operator, \(P\) is a coherent-sector projector,
\(\chi(A)\) is a spectral window, and
\[
  \ta
  =
  \lambda_\ast^{-1}\log(C_Q/\epsadm)
\]
is the earliest proper-time/heat-time certified by the stated exponential leakage bound. It
is not automatically the exact first admissible time or a first-principles physical scale.
In the local representation, the kernel
\[
  \Kadm(x,y)=\langle x|\Badm|y\rangle
\]
defines a finite localization profile and recovers the Dirac delta only in the sharp limit.
This is the particle shadow. In the spectral representation,
\[
  \Kadm(x,y)
  =
  \sum_{n\in \coh}
  \chi(\lambda_n)^2 e^{-\ta\lambda_n}
  \phi_n(x)\phi_n^\ast(y),
\]
the same object appears as a weighted coherent modal structure whose phase evolution
produces interference. This is the wave shadow.

The paper organizes this conditional encoding in three ways. First, the delta limit is stated as
a distributional convergence theorem. Second, branch damping is formulated as a Schur
product quantum channel, with the positivity condition \(D\succeq0\) ensuring complete
positivity and trace preservation. Third, finite detector effects are treated as normalized
POVM kernels and are completed by quantum instruments when outcome probabilities and
post-measurement states are discussed. A survivor-basin mechanism is an additional MTT
interpretation, not a consequence of the POVM axioms. The double-slit experiment is a
canonical example:
branch coherence gives interference, while which-way selection damps off-diagonal branch
terms and leaves localized events.

The rigorous result is an operator dual-representation theorem: one positive filtered operator
has local and spectral representations, while valid POVMs, instruments, and Schur channels
model localization and coherence loss. This reproduces standard quantum predictions only
because the quantum state, dynamics, effects, and instrument are supplied as inputs. The
claim that all of them descend from one MTT source remains a conditional encoding theorem,
and the Born rule requires an independent basin-measure intertwiner.
\end{abstract}

\section*{Version 5 Revision Note}
\begin{description}
\item[Supersedes] The April 2026 Version 4 manuscript with the same subject.
\item[Reason] The earlier paper contained valid kernel, POVM, and Schur-channel results but
could be read as deriving measurement outcomes, the Born rule, and standard quantum
mechanics from the existence of two kernel representations.
\item[Resolution] This version distinguishes exact and certified damping times, adds the
quantum-instrument layer required for operational measurement, states the exact
Born-bridge condition, and classifies projection duality as a conditional MTT encoding.
\item[Retained result] The local/spectral kernel identity, distributional sharp limit,
positive Schur-channel criterion, and visibility calculations are retained.
\item[Remaining boundary] A selected common MTT source must derive the state, Hamiltonian,
effects, instrument, branch kernel, and basin measure without importing their measured
predictions.
\end{description}

\tableofcontents

\section{Purpose and claim discipline}
\label{sec:purpose}

Wave--particle duality is one of the oldest signs that effective physical description is
projection-dependent. In one experimental arrangement a microscopic system is described by
a phase-coherent amplitude that propagates, diffracts, and interferes. In another arrangement
the same system produces a localized detector record. Standard quantum mechanics encodes
both facts successfully, but it normally treats their coexistence as a primitive feature of the
formalism: states evolve as amplitudes, while measurements return localized outcomes.

The purpose of this paper is to show that this dual appearance follows naturally from the
projection architecture of Modal Triplet Theory (MTT). The preceding delta--projection
sequence established that Dirac delta distributions should not be read as primitive physical
objects whenever finite coherent support is relevant. Rather, they arise as singular downstream
limits of finite admissible kernels. The fixed-point kernel construction then replaced arbitrary
smoothing by the canonical admissible operator
\begin{equation}
  \Badm
  =
  P\chi(A)e^{-\ta A}\chi(A)P,
  \label{eq:purpose-badm}
\end{equation}
where \(A\) is the declared fixed-point stabilization operator, \(P\) is the coherent-sector
projector, \(\chi(A)\) is the spectral window, and \(\ta\) is the proper-time/heat-time
certified by the declared damping bound.

The present paper applies this same object to the wave--particle problem. The claim is not
that there are two underlying things, one wave-like and one particle-like. Nor is the claim
that an object changes its ontology when a detector is introduced. The claim is that a single
finite coherent-sector excitation has two different downstream representations. In a local
or source representation, the admissible kernel appears as a finite localization profile whose
sharp idealization is a Dirac delta. In a spectral or phase representation, the same kernel
appears as a weighted coherent modal structure whose retained phases generate interference.

\begin{quote}
\emph{Central claim.}
Wave-like interference and particle-like localization are two downstream projection shadows
of one finite coherent-sector excitation.
\end{quote}

Equivalently,
\begin{equation}
  \boxed{
  \text{wave--particle duality}
  \quad\leadsto\quad
  \text{a conditional projection-duality encoding}.
  }
  \label{eq:purpose-projection-duality}
\end{equation}

The word ``shadow'' is used in the technical MTT sense. A shadow is not an illusion, nor is it
a merely subjective description. It is a stable effective representation obtained after projecting
a richer coherent structure into a lower-capacity descriptive context. The same upstream
coherent excitation may therefore have inequivalent downstream shadows, depending on which
projection context is imposed.

\subsection{What is being claimed}
\label{subsec:claimed}

This paper proves and organizes the following structural claims.

First, the admissible coherent kernel has both a local representation and a spectral
representation. In local notation,
\begin{equation}
  \Kadm(x,y)=\langle x|\Badm|y\rangle
\end{equation}
acts as a finite source, response, or detector-amplitude kernel. In spectral notation, under
the assumptions stated below,
\begin{equation}
  \Kadm(x,y)
  =
  \sum_{n\in\coh}
  \chi(\lambda_n)^2 e^{-\ta\lambda_n}
  \phi_n(x)\phi_n^\ast(y).
  \label{eq:purpose-spectral-kernel}
\end{equation}
The local representation is the particle-like side. The spectral representation is the
wave-like side.

Second, the point-particle description is recovered only as a limiting representation. In a
sharp spectral/proper-time limit,
\begin{equation}
  \Kadm(x,y)\longrightarrow \delta(x-y)
\end{equation}
in the distributional sense. Thus MTT does not postulate literal point particles. It explains
why point-particle descriptions are successful when the finite coherent width is below the
resolution scale of the probing apparatus.

Third, wave-like behavior is identified operationally with preserved off-diagonal coherence:
superposition, retained relative phase, and observable cross terms. If two branches \(a,b\)
remain in the same coherent sector, then interference terms survive. If a measurement
arrangement distinguishes the branches strongly enough, those off-diagonal terms are damped.

Fourth, branch damping is not introduced as an arbitrary visibility factor. It must define a
valid quantum channel on the branch density matrix. In finite branch bases this means that
the damping matrix \(D=(D_{ab})\) must satisfy
\begin{equation}
  D\succeq0,
  \qquad
  D_{aa}=1,
\end{equation}
so that the Schur map
\begin{equation}
  \rho\mapsto D\circ\rho
\end{equation}
is completely positive and trace preserving. This condition is essential: not every proposed
interference-damping rule is physically admissible.

Fifth, measurement is represented first by standard operational data and only then compared
with a possible survivor-basin completion. A detector context \(\Meas\) has device-sector data
\begin{equation}
  (A_{\Meas},P_{\Meas},\chi_{\Meas},\tau_{\Meas},
  \{E_i^{(\Meas)}\}_{i\in I},
  \{\mathcal I_i^{(\Meas)}\}_{i\in I},
  \mathfrak B_{\Meas}).
\end{equation}
Here the effects and instrument define probabilities and state updates;
\(\mathfrak B_{\Meas}=\{B_i^{(\Meas)}\}_{i\in I}\) is an additional proposed survivor-basin partition.
Different devices can therefore disturb the same coherent excitation in different ways, select
different branch bases, and induce different damping matrices
\begin{equation}
  D_{ab}^{(\Meas)}
  =
  \exp[-\tau_{\Meas}\Lambda_{ab}^{(\Meas)}],
\end{equation}
when the exponent defines an admissible positive kernel. Identifying an instrument outcome
with a stabilized basin record requires a further intertwining theorem; it does not follow
from positivity of the effects or damping channel alone.

\subsection{What is not being claimed}
\label{subsec:not-claimed}

This paper does not claim that standard quantum mechanics is false. It does not deny the
usefulness of wavefunctions, Hilbert spaces, amplitudes, density matrices, POVMs, or
projection-valued idealizations. Instead, it explains why these structures appear as effective
descriptions inside finite admissible regimes.

This paper also does not claim that particles are ordinary extended classical blobs. The MTT
claim is more precise. The upstream object is not a little ball, and it is not a classical wave
spread through space. It is a stable coherent-sector excitation. That excitation can appear
pointlike when probed locally and wavelike when represented spectrally.

Thus
\begin{equation}
  \text{particle}\neq \text{literal mathematical point},
\end{equation}
and
\begin{equation}
  \text{wave}\neq \text{separate physical substance}.
\end{equation}
Rather,
\begin{equation}
  \text{particle shadow}
  =
  \text{localized delta-limit representation},
\end{equation}
while
\begin{equation}
  \text{wave shadow}
  =
  \text{spectral phase-coherent representation}.
\end{equation}
Both are generated by the same admissible MTT structure.

Nor does this paper compute all sector-specific numerical parameters. The quantities
\(A\), \(P\), \(\chi\), \(\lambda_\ast\), \(C_Q\), \(\epsadm\), and the relevant detector or
branch-separation metric must be supplied by the physical sector under consideration. The
present claim is conditional but sharp: once those data are supplied by an admissible
fixed-point regime, no independent wave--particle switch and no arbitrary smoothing width
remain.

\subsection{Relation to probability and the Born rule}
\label{subsec:born-relation}

A common source of confusion is to conflate two different questions:
\begin{align}
  &\text{Which outcome is selected?}
  \label{eq:which-outcome}\\
  &\text{How much interference remains before or during selection?}
  \label{eq:how-much-interference}
\end{align}
MTT treats these as related but distinct questions.

The first question is answered operationally by the supplied POVM or instrument:
\[
  p_i=\tr(\rho E_i)=\tr\mathcal I_i(\rho).
\]
The Born-rule shadow-bridge program proposes to reproduce these probabilities by relative
measures of stabilized selection basins:
\begin{equation}
  \mathbb P(i)
  =
  \frac{\mu(B_i)}{\sum_j\mu(B_j)},
  \label{eq:basin-born}
\end{equation}
where \(B_i\) is the basin of the \(i\)-th stabilized outcome and \(\mu\) is the relevant
admissibility-weighted basin measure. Exact agreement requires the bridge condition
\[
  \frac{\mu_\rho(B_i)}{\sum_j\mu_\rho(B_j)}
  =
  \tr(\rho E_i)
  \qquad\text{for every admitted }\rho\text{ and }i.
\]
This condition is not proved by the existence of the two kernel representations.

The second question is answered by branch-coherence survival. Before a unique outcome is
stabilized, or in arrangements where alternatives are recombined, the off-diagonal terms
\(\rho_{ab}\) determine interference visibility. Device-induced disturbance acts on these
terms by an admissible damping channel,
\begin{equation}
  \rho_{ab}\longmapsto D_{ab}^{(\Meas)}\rho_{ab}.
  \label{eq:branch-damping-purpose}
\end{equation}
This is the visibility side of measurement.

Thus this paper does not replace Born probabilities with visibility damping. Instead it keeps
the two roles separate:
\begin{equation}
  \boxed{
  \text{a proved basin--Born bridge would reproduce stabilized outcome frequencies},
  }
\end{equation}
whereas
\begin{equation}
  \boxed{
  \text{branch coherence determines interference visibility}.
  }
\end{equation}
This distinction is crucial for the double-slit experiment. The screen distribution is shaped
by coherent interference when \(D_{12}\approx1\), while individual screen records are still
localized detector outcomes whose frequencies are governed by the relevant measurement
effects and basin weights.

\subsection{Relation to measurement as disturbance and stabilization}
\label{subsec:disturbance-relation}

The measurement picture used here is the same one used throughout the MTT measurement
sequence. A measuring device is not a passive reader of a pre-existing classical value. It is a
localized disturbance coupled to a finite-capacity coherent sector. After the disturbance, the
system undergoes re-coherence: high-frequency or inadmissible components are damped, the
coherent projector selects the retained content, and the resulting state stabilizes into one of
the available survivor basins.

In symbolic form, a measurement context induces the sequence
\begin{equation}
  \begin{aligned}
  \text{coherent evolution}
  &\longrightarrow \text{localized disturbance}
  \longrightarrow \text{projection}\\
  &\longrightarrow \text{survivor-basin stabilization}.
  \end{aligned}
  \label{eq:evolve-project-stabilize-purpose}
\end{equation}
Different devices implement different versions of this sequence. A weak which-way marker,
a strong absorbing detector, a Stern--Gerlach magnet, a delayed-choice recombination stage,
and a phase-sensitive interferometer do not impose the same disturbance. They define different
measurement contexts \(\Meas\), different branch partitions, and different damping kernels.

This is why the device matters in the double-slit experiment. With no which-way device,
the two path branches remain mutually coherent and the off-diagonal term survives. With a
weak which-way device, the off-diagonal term is partially damped. With a strong which-way
device, the path branches are separated and the interference term is suppressed. The
instrument supplies a definite operational record. Calling that record a stabilized survivor
basin additionally invokes the declared instrument--basin intertwiner.

\subsection{Why this is not merely decoherence}
\label{subsec:not-merely-decoherence}

The formal appearance of a damping factor
\begin{equation}
  \rho_{ab}\mapsto D_{ab}\rho_{ab}
\end{equation}
resembles ordinary decoherence. The relation is real but incomplete. Environmental
decoherence explains how entanglement with uncontrolled degrees of freedom suppresses
off-diagonal terms in a reduced density matrix. MTT includes such suppression as one possible
shadow, but the present structure is broader.

In MTT, branch damping is tied to admissibility, coherent projection, detector-sector
disturbance, and survivor-basin stabilization. The damping factor is not merely a phenomenological
loss of phase coherence. It must be compatible with a finite measurement kernel, a valid
quantum channel, and a basin-selection process. Decoherence suppresses interference; MTT
also asks which stabilized branch is selected, which basin measure governs its frequency, and
which finite coherent kernel underlies the local and spectral shadows in the first place.

\subsection{Paper structure}
\label{subsec:paper-structure}

The paper proceeds as follows. The next section fixes the analytic setting and the admissible
kernel. We then state the precise delta-limit theorem, giving the mathematical basis for the
particle shadow. The following section develops the spectral representation and coherent
phase evolution, giving the mathematical basis for the wave shadow. We then formulate
finite measurement effects, Schur-channel branch damping, and the operator version of the
double-slit experiment. After that we discuss device-dependent disturbance, further
experimental modules, recovery of standard quantum mechanics, relation to existing accounts,
and the failure modes of the construction.

The intended result is a rigorous structural synthesis:
\begin{equation}
  \boxed{
  \text{one admissible coherent kernel}
  \quad\Longrightarrow\quad
  \begin{cases}
  \text{delta-like local particle shadow},\\
  \text{spectral phase wave shadow},\\
  \text{finite measurement-selection shadow}.
  \end{cases}
  }
  \label{eq:purpose-summary}
\end{equation}


\section{Analytic setting and admissible coherent kernels}
\label{sec:analytic-setting}

This section fixes the analytic assumptions used in the structural theorems below. The goal is
not to introduce a new quantum formalism, but to identify the minimal operator-theoretic data
needed for the local and spectral shadows of a coherent MTT excitation to be well-defined.

\subsection{Base Hilbert space and fixed-point operator}
\label{subsec:base-hilbert-space}

Let \(X\) be a compact smooth Riemannian manifold without boundary, or a compact manifold
with boundary equipped with a self-adjoint elliptic boundary condition. More generally, \(X\)
may be replaced by a compact admissible chart and scalar functions by sections of a Hermitian
vector bundle. We write \(L^2(X)\) for the Hilbert space of square-integrable sections.

Let
\begin{equation}
  A\ge 0
\end{equation}
be a nonnegative self-adjoint elliptic operator on \(L^2(X)\), with compact resolvent. The
spectral theorem gives an orthonormal basis of smooth eigenfunctions or eigensections
\begin{equation}
  A\phi_n=\lambda_n\phi_n,
  \qquad
  0\le \lambda_0\le\lambda_1\le\cdots,
  \label{eq:A-spectrum}
\end{equation}
where eigenvalues are repeated according to multiplicity. In MTT, \(A\) is interpreted as the
linearized fixed-point stabilization, damping, or admissibility operator in the coherent chart
under consideration.

The compact setting is used to state the spectral results without unnecessary measure-theoretic
technicalities. In flat noncompact charts, such as \(X=\mathbb R^d\), the corresponding statements
are obtained by Fourier transform and convergence on Schwartz test functions. For example,
\(A=-\Delta\) has continuous spectral parameter \(k^2\), and heat kernels are interpreted as
tempered-distribution kernels.

\begin{assumption}[Fixed-point spectral data]
\label{ass:fixed-point-spectral-data}
The fixed-point coherent sector is specified by:
\begin{enumerate}[label=(\roman*)]
  \item a nonnegative self-adjoint elliptic operator \(A\) with compact resolvent;
  \item an isolated coherent spectral cluster \(\Omega_{\rm coh}\subset\sigma(A)\);
  \item the Riesz spectral projector
  \begin{equation}
    P
    =
    \frac{1}{2\pi i}
    \oint_\Gamma (z-A)^{-1}\,\dd z,
    \label{eq:riesz-P}
  \end{equation}
  where \(\Gamma\) encloses \(\Omega_{\rm coh}\) and no discarded spectrum;
  \item the complementary projector
  \begin{equation}
    Q=I-P;
  \end{equation}
  \item a discarded-sector gap
  \begin{equation}
    A|_{\Ran Q}\ge \lambda_\ast I,
    \qquad
    \lambda_\ast>0;
    \label{eq:discarded-gap}
  \end{equation}
  \item an admissible spectral window \(\chi:[0,\infty)\to[0,1]\) subordinate to the coherent
  spectral split.
\end{enumerate}
\end{assumption}

Because \(P\) is a Riesz projector of \(A\), it commutes with \(A\) and preserves the smooth
spectral domain. This avoids a common ambiguity: a merely bounded projector commuting with
\(A\) is not enough, by itself, to guarantee the smooth kernel properties used below. In the
sharp cluster realization, \(\chi(A)=P\). Smooth windows may be used when a transition layer
inside a gap-protected admissible neighborhood is required, but such windows are subordinate
to the same spectral split and are not independent phenomenological regulators.

\begin{definition}[Coherent and discarded sectors]
\label{def:coh-discarded-sectors}
The retained coherent sector and discarded sector are
\begin{equation}
  \calH_{\rm coh}:=\Ran P,
  \qquad
  \calH_{\rm inc}:=\Ran Q.
\end{equation}
A vector \(\Psi\in L^2(X)\) decomposes as
\begin{equation}
  \Psi=P\Psi+Q\Psi.
\end{equation}
The coherent component \(P\Psi\) is the part represented by the effective physical description.
The discarded component \(Q\Psi\) is not necessarily nonexistent; it is inadmissible for the
chosen effective chart.
\end{definition}

\subsection{Exact and certified proper-time scales}
\label{subsec:damping-selected-time}

The filtered kernel contains a heat/proper-time parameter. An exponential damping estimate
can certify an admissible time, but a bound need not identify the exact first crossing, and
neither quantity is physical until its operator and tolerance are selected.

Let \(\Phi_t\) denote the local linearized or locally controlled damping flow. In the simplest
self-adjoint linear case,
\begin{equation}
  \Phi_t=e^{-tA}.
\end{equation}
Assume that the discarded sector satisfies
\begin{equation}
  \norm{Q\Phi_tQ}
  \le
  C_Qe^{-\lambda_\ast t},
  \qquad
  t\ge0,
  \label{eq:discarded-sector-bound}
\end{equation}
for some \(C_Q\ge1\). In the canonical self-adjoint spectral case with
\(A|_{\Ran Q}\ge\lambda_\ast I\), one has \(C_Q=1\).

\begin{definition}[Admissible residual leakage]
\label{def:admissible-residual-leakage}
Let \(0<\epsadm<1\). A stabilization step of heat/proper-time \(t\) is
\(\epsadm\)-admissible if
\begin{equation}
  \norm{Q\Phi_tQ}\le \epsadm.
\end{equation}
The parameter \(\epsadm\) is a regime-dependent admissibility tolerance. It may be determined
by basin separation, detector resolution, closure strain, or other sector data. It is not a
universal constant in this paper.
\end{definition}

\begin{definition}[Exact exit time]
\label{def:exact-exit-time}
Define
\begin{equation}
  \tau_{\rm exit}
  :=
  \inf\{t\ge0:\norm{Q\Phi_tQ}\le\epsadm\}.
\end{equation}
\end{definition}

\begin{proposition}[Certified admissible heat/proper-time]
\label{prop:minimal-admissible-time}
Under the damping estimate \eqref{eq:discarded-sector-bound}, every
\begin{equation}
  t\ge
  \lambda_\ast^{-1}
  \log\frac{C_Q}{\epsadm}
\end{equation}
is \(\epsadm\)-admissible. The earliest time certified by this bound is
\begin{equation}
  \ta
  :=
  \lambda_\ast^{-1}
  \log\frac{C_Q}{\epsadm}.
  \label{eq:tau-adm}
\end{equation}
If the leakage norm is nonincreasing, then
\begin{equation}
  \tau_{\rm exit}\le\ta.
\end{equation}
Equality requires additional information. In the reducing self-adjoint case it follows from
the exact spectral norm formula, as proved in the dedicated admissibility-time paper.
\end{proposition}

\begin{proof}
The bound \eqref{eq:discarded-sector-bound} gives
\begin{equation}
  \norm{Q\Phi_tQ}
  \le
  C_Qe^{-\lambda_\ast t}.
\end{equation}
Requiring the right-hand side to be at most \(\epsadm\) gives
\begin{equation}
  C_Qe^{-\lambda_\ast t}\le\epsadm.
\end{equation}
Taking logarithms yields
\begin{equation}
  t\ge
  \lambda_\ast^{-1}\log(C_Q/\epsadm).
\end{equation}
The right-hand side is an upper bound for the actual leakage, so its crossing time may be
later than the exact exit time.
\end{proof}

\begin{remark}[Heat-time rather than clock time]
\label{rem:heat-time}
The parameter \(\ta\) is a proper-time/heat-time parameter for the semigroup generated by
\(A\). If \(A\) is Laplace-type and compared to \(k^2\), then in units \(\hbar=c=1\),
\begin{equation}
  [A]=E^2,
  \qquad
  [\ta]=E^{-2}.
\end{equation}
The corresponding coherence length scale is
\begin{equation}
  \ell_{\rm coh}\sim \sqrt{\ta},
\end{equation}
and the effective coherence-energy scale is
\begin{equation}
  \Lambda_{\rm eff}\sim \ta^{-1/2}.
\end{equation}
\end{remark}

\subsection{The admissible coherent kernel}
\label{subsec:admissible-kernel}

\begin{definition}[Conditionally specified coherent operator]
\label{def:admissible-coherent-operator}
Given the declared data of Assumption~\ref{ass:fixed-point-spectral-data} and the certified
scale \(\ta\), define
\begin{equation}
  \Badm
  :=
  P\chi(A)e^{-\ta A}\chi(A)P.
  \label{eq:Badm-definition}
\end{equation}
When \(\Badm\) has an integral kernel, we write
\begin{equation}
  \Kadm(x,y)
  :=
  \langle x|\Badm|y\rangle.
  \label{eq:Kadm-definition}
\end{equation}
\end{definition}

Because \(A\) is self-adjoint and \(P,\chi(A)\) are defined by functional calculus, the operator
\(\Badm\) is positive and bounded. For \(\ta>0\), elliptic heat-kernel regularity implies that
\(\Badm\) is smoothing on the retained spectral sector. In a finite-rank coherent cluster,
\(\Badm\) has a smooth kernel even at \(\ta=0\).

\begin{proposition}[Spectral form of the admissible kernel]
\label{prop:spectral-form}
Under Assumption~\ref{ass:fixed-point-spectral-data}, suppose \(A\phi_n=\lambda_n\phi_n\) is an
orthonormal spectral resolution. Then
\begin{equation}
  \Badm f
  =
  \sum_n
  p_n\chi(\lambda_n)^2 e^{-\ta\lambda_n}
  \langle \phi_n,f\rangle\phi_n,
  \label{eq:Badm-spectral-action}
\end{equation}
where \(p_n=1\) if \(\phi_n\in\Ran P\) and \(p_n=0\) otherwise. Consequently,
\begin{equation}
  \Kadm(x,y)
  =
  \sum_{n\in\coh}
  \chi(\lambda_n)^2 e^{-\ta\lambda_n}
  \phi_n(x)\phi_n^\ast(y).
  \label{eq:Kadm-spectral}
\end{equation}
\end{proposition}

\begin{proof}
Since \(P\), \(\chi(A)\), and \(e^{-\ta A}\) are all functions of \(A\) on the spectral
decomposition, they are diagonal in the same eigenbasis. Applying the product
\(P\chi(A)e^{-\ta A}\chi(A)P\) to \(f=\sum_n\langle\phi_n,f\rangle\phi_n\) gives
\eqref{eq:Badm-spectral-action}. The kernel formula follows by writing the corresponding
Schwartz kernel of the spectral expansion.
\end{proof}

\begin{remark}[Amplitude kernel versus probability density]
\label{rem:kernel-not-density}
The kernel \(\Kadm(x,y)\) should not be assumed to be a probability density. In general it is
a finite source, response, or amplitude kernel and may have oscillatory structure depending on
the spectral window. Positive probabilities are obtained from positive effects, for example
\begin{equation}
  E_x=|k_x\rangle\langle k_x|,
\end{equation}
where \(k_x(y)=K_{\rm det}(y,x)\), or more generally from a POVM satisfying the normalization
conditions stated below. This distinction prevents the pointlike ``particle shadow'' from being
misread as a classical density blob.
\end{remark}

\subsection{Local and spectral readings}
\label{subsec:local-spectral-readings}

The same kernel \(\Kadm\) has two readings.

First, fixing \(y=x_0\), the function
\begin{equation}
  x\longmapsto \Kadm(x,x_0)
\end{equation}
is a finite local response to a source centered at \(x_0\). If the effective kernel width is below
the resolving power of a detector, this response is operationally indistinguishable from a point
source. This is the local or particle-like shadow.

Second, the expansion \eqref{eq:Kadm-spectral} expresses the same kernel as a weighted sum
of coherent modes. If a coherent phase generator \(H_{\rm coh}\) transports these retained
modes unitarily, then relative phases between modal or branch components produce interference.
This is the spectral or wave-like shadow.

The paper's central thesis is that these are not two objects. They are two representations of
one admissible coherent operator.

\subsection{Detector-sector kernels}
\label{subsec:detector-sector-kernels}

Measurement introduces device-specific data. A detector is not modeled by an arbitrary
external delta function. It is represented by a detector-sector instance of the same admissible
kernel construction.

\begin{definition}[Detector-sector admissible kernel]
\label{def:detector-sector-kernel}
A measurement context \(\Meas\) supplies detector-sector data
\begin{equation}
  (A_{\Meas},P_{\Meas},\chi_{\Meas},\tau_{\Meas}).
\end{equation}
The associated detector filter is denoted
\begin{equation}
  F_{\Meas}
  =
  P_{\Meas}\chi_{\Meas}(A_{\Meas})
  e^{-\tau_{\Meas}A_{\Meas}}
  \chi_{\Meas}(A_{\Meas})P_{\Meas}.
  \label{eq:detector-filter-F}
\end{equation}
Its kernel is denoted
\begin{equation}
  K_{\Meas}(y,x)
  =
  \langle y|F_{\Meas}|x\rangle.
  \label{eq:detector-kernel-F}
\end{equation}
When the context is a position-sensitive detector we also write
\begin{equation}
  K_{\rm det}(y,x):=K_{\Meas}(y,x).
\end{equation}
\end{definition}

Thus \(K_{\rm det}\) is a detector-sector realization of the same proposed kernel
architecture. This definition does not derive it from the system kernel. A physical model
must supply a connection-preserving source map between the two sectors. Different devices may have
different \(A_{\Meas}\), \(P_{\Meas}\), \(\chi_{\Meas}\), and \(\tau_{\Meas}\), and therefore disturb and
select the coherent excitation differently.

\subsection{Finite effects and normalized detection}
\label{subsec:finite-effects}

Given a detector kernel \(K_{\rm det}\), define the rank-one effect density
\begin{equation}
  E_x(y,z)
  =
  K_{\rm det}(y,x)K_{\rm det}^\ast(z,x).
  \label{eq:finite-effect-kernel}
\end{equation}
Equivalently,
\begin{equation}
  E_x=|k_x\rangle\langle k_x|,
  \qquad
  k_x(y)=K_{\rm det}(y,x).
\end{equation}

For a probability-conserving detector on the retained sector, the effects satisfy the POVM
normalization
\begin{equation}
  \int_X E_x\,\dd x
  =
  I_{\rm coh}.
  \label{eq:povm-normalization}
\end{equation}
Then a state \(\rho\) gives detection density
\begin{equation}
  p(x)=\tr(\rho E_x),
\end{equation}
and
\begin{equation}
  \int_X p(x)\,\dd x
  =
  \tr(\rho I_{\rm coh}).
\end{equation}
For normalized states on the retained sector, this equals one.

\begin{definition}[Instrument completion]
\label{def:instrument-completion}
A POVM specifies outcome probabilities but not the state update. An operational
measurement therefore also requires a quantum instrument
\(\{\mathcal I_i\}_{i\in I}\): each \(\mathcal I_i\) is completely positive and
trace-nonincreasing, \(\sum_i\mathcal I_i\) is trace preserving for a closed detector, and
\[
  E_i=\mathcal I_i^\ast(I_{\rm coh}).
\]
For input state \(\rho\),
\[
  p_i=\tr\mathcal I_i(\rho)=\tr(\rho E_i),
  \qquad
  \rho_i'=\frac{\mathcal I_i(\rho)}{p_i}
\]
when \(p_i>0\).
\end{definition}

\begin{remark}[What the POVM does not prove]
The effects \(\{E_i\}\) determine the Born probabilities only after the quantum state and
Born pairing are supplied. They do not select a unique instrument, do not prove a unique
outcome, and do not identify an MTT survivor basin. Those are separate promotion steps.
\end{remark}

\begin{remark}[Open detectors and leakage]
\label{rem:open-detectors}
If the detector has unobserved discarded channels, then one may instead have
\begin{equation}
  \int_X E_x\,\dd x
  \le
  I_{\rm coh}.
\end{equation}
The missing probability corresponds to loss into channels not recorded by the chosen detector
context. The closed-detector case \eqref{eq:povm-normalization} is used when discussing
ordinary normalized detection statistics.
\end{remark}

\subsection{Operational meanings of particle-like and wave-like}
\label{subsec:operational-meanings}

The words ``particle-like'' and ``wave-like'' are used operationally.

\begin{definition}[Particle-like behavior]
\label{def:particle-like}
An excitation is particle-like relative to a detector context \(\Meas\) when:
\begin{enumerate}[label=(\roman*)]
  \item detector effects \(E_i^{(\Meas)}\) stabilize discrete records;
  \item the relevant detector kernel width is below the resolution scale of the apparatus;
  \item the sharp idealization of those effects is a delta/PVM limit;
  \item repeated trials yield localized records distributed according to the appropriate
  measurement probabilities.
\end{enumerate}
\end{definition}

\begin{definition}[Wave-like behavior]
\label{def:wave-like}
An excitation is wave-like relative to an experimental context when:
\begin{enumerate}[label=(\roman*)]
  \item it supports coherent superposition of retained branches or modes;
  \item relative phases between those branches or modes are preserved;
  \item observable probabilities contain off-diagonal cross terms;
  \item those cross terms are visible as interference, diffraction, or phase-sensitive
  modulation.
\end{enumerate}
\end{definition}

These definitions make the main claim precise. MTT does not assert that an entity is
``really'' a wave or ``really'' a particle. It asserts that particle-like and wave-like behavior
are context-dependent downstream manifestations of a coherent-sector excitation.



\section{The particle shadow: localization as a delta limit}
\label{sec:particle-shadow}

This section makes precise the sense in which pointlike particle behavior is a limiting
description of finite coherent support. The claim is not that MTT replaces particles by
classical extended objects. The claim is that the mathematical point source is the sharp
idealization of a finite admissible kernel.

\subsection{Finite local response}
\label{subsec:finite-local-response}

Let \(x_0\in X\). The local response of the coherent sector to a source centered at \(x_0\)
is represented by the kernel section
\begin{equation}
  x\longmapsto \Kadm(x,x_0).
  \label{eq:local-response-section}
\end{equation}
In the ideal point-particle description one instead writes
\begin{equation}
  x\longmapsto \delta(x-x_0).
\end{equation}
MTT reverses the hierarchy. The finite object \(\Kadm(x,x_0)\) is native to a finite
admissible regime; the Dirac delta is recovered only after an idealizing limit.

For a source amplitude \(J\) localized near \(x_0\), the coherent response is
\begin{equation}
  (B_{\rm adm}J)(x)
  =
  \int_X \Kadm(x,y)J(y)\,\dd y.
  \label{eq:kernel-response-source}
\end{equation}
If \(J=\delta_{x_0}\) is used as an ideal test distribution, then
\begin{equation}
  (B_{\rm adm}\delta_{x_0})(x)
  =
  \Kadm(x,x_0).
\end{equation}
Thus even a formally point-supported input is represented downstream by a finite coherent
response.

\begin{remark}[Not a classical matter density]
\label{rem:not-classical-density}
The function \(x\mapsto\Kadm(x,x_0)\) is a response amplitude or kernel section. It should
not be read as a positive classical matter density unless additional positivity conditions are
present. Detection probabilities are computed from positive effects, such as
\(E_x=|k_x\rangle\langle k_x|\), or from a normalized POVM. This distinction matters:
the particle shadow is a localization limit of the coherent kernel, not a claim that particles
are small classical blobs.
\end{remark}

\subsection{Resolution-dependent pointlikeness}
\label{subsec:resolution-dependent-pointlikeness}

Let \(\ell_{\rm coh}\) denote the characteristic width of the finite coherent kernel in the local
chart. In heat-kernel models on a flat chart,
\begin{equation}
  \ell_{\rm coh}\sim \sqrt{\ta}.
  \label{eq:ell-coh}
\end{equation}
More precisely, for the Euclidean heat kernel
\begin{equation}
  K_\tau(x,y)
  =
  (4\pi\tau)^{-d/2}
  \exp\left(-\frac{|x-y|^2}{4\tau}\right),
  \label{eq:flat-heat-kernel}
\end{equation}
the root-mean-square width is of order \(\sqrt{\tau}\), with convention-dependent numerical
factors.

Let \(\ell_{\rm res}\) be the spatial resolution of the probing apparatus. The coherent excitation
is pointlike relative to that apparatus when
\begin{equation}
  \ell_{\rm coh}\ll \ell_{\rm res}.
  \label{eq:pointlike-resolution}
\end{equation}
Equivalently, if the experiment probes momenta up to \(E_{\rm probe}\), while the coherent
scale is
\begin{equation}
  \Lambda_{\rm eff}\sim \ta^{-1/2},
\end{equation}
then the pointlike approximation is valid when
\begin{equation}
  E_{\rm probe}\ll \Lambda_{\rm eff}.
  \label{eq:pointlike-energy}
\end{equation}

\begin{definition}[Effective point particle]
\label{def:effective-point-particle}
An excitation is an effective point particle relative to a detector context \(\Meas\) if its
detector-sector localization width is below the resolution scale of \(\Meas\), so that the
finite effects \(E_x^{(\Meas)}\) are operationally indistinguishable from sharp position
effects on the retained sector.
\end{definition}

Thus pointlikeness is not absolute. It is a relation between a coherent support scale and a
probing scale.

\subsection{Sharp spectral/proper-time limit}
\label{subsec:sharp-limit}

We now state the precise delta-limit theorem used throughout the paper. The statement is
standard spectral analysis, but its interpretation is central for MTT.

Let \(P_N=\mathbf 1_{[0,\Lambda_N]}(A)\) be the spectral projector of \(A\) onto eigenvalues
at most \(\Lambda_N\), with
\begin{equation}
  \Lambda_N\to\infty.
\end{equation}
Let
\begin{equation}
  B_{N,\tau}
  :=
  P_N e^{-\tau A}P_N.
  \label{eq:BNtau}
\end{equation}
When \(B_{N,\tau}\) has a kernel, denote it by \(K_{N,\tau}(x,y)\).

\begin{theorem}[Delta limit of finite coherent kernels]
\label{thm:delta-limit}
Let \(X\) be compact and \(A\ge0\) be self-adjoint elliptic with compact resolvent. Let
\(\tau_N\downarrow0\) and \(\Lambda_N\to\infty\). Then, for every \(f\in C^\infty(X)\),
\begin{equation}
  B_{N,\tau_N}f\longrightarrow f
\end{equation}
in \(C^\infty(X)\). Equivalently,
\begin{equation}
  K_{N,\tau_N}(x,y)
  \longrightarrow
  \delta(x-y)
\end{equation}
in the distributional sense on \(X\times X\).
\end{theorem}

\begin{proof}
Write
\begin{equation}
  f-B_{N,\tau_N}f
  =
  (I-P_N)f
  +
  P_N(I-e^{-\tau_NA})P_N f.
  \label{eq:delta-proof-split}
\end{equation}
The first term tends to zero in all Sobolev norms because spectral projectors converge
strongly to the identity on smooth functions, and smooth functions have rapidly decaying
spectral coefficients.

For the second term, fix a Sobolev index \(s\). Choose \(M>s\). Since \(f\in C^\infty(X)\),
\begin{equation}
  \sum_n (1+\lambda_n)^M |\langle\phi_n,f\rangle|^2<\infty.
\end{equation}
For each fixed \(n\),
\begin{equation}
  1-e^{-\tau_N\lambda_n}\to0.
\end{equation}
Moreover, the Sobolev-weighted summands are dominated by a summable multiple of
\((1+\lambda_n)^M|\langle\phi_n,f\rangle|^2\). Dominated convergence gives convergence
to zero in \(H^s\). Since \(s\) is arbitrary, Sobolev embedding gives convergence in
\(C^\infty(X)\). The distributional convergence of kernels is the kernel formulation of
strong convergence to the identity on test functions.
\end{proof}

\begin{remark}[Meaning of the sharp limit]
\label{rem:meaning-sharp-limit}
The theorem does not say that the physical coherent kernel is literally a Dirac delta. It says
that the Dirac delta is recovered when:
\begin{enumerate}[label=(\roman*)]
  \item coherent bandwidth becomes complete;
  \item the proper-time width tends to zero;
  \item no discarded sector remains unresolved.
\end{enumerate}
A finite-capacity coherent regime does not satisfy these idealizations exactly. Its native
object is the finite kernel.
\end{remark}

\subsection{Finite-width correction estimates}
\label{subsec:finite-width-estimates}

The sharp limit is useful, but finite-width estimates are more physically meaningful. The
difference between a sharp retained description and the admissibly filtered description has
three conceptually distinct sources:
\begin{enumerate}[label=(\roman*)]
  \item loss from discarding the noncoherent sector;
  \item loss from the spectral acceptance window \(\chi(A)\);
  \item smoothing from the finite heat/proper-time factor \(e^{-\tau_{\rm adm}A}\).
\end{enumerate}

For \(f\in L^2(X)\), using
\begin{equation}
  B_{\rm adm}=P\chi(A)e^{-\tau_{\rm adm}A}\chi(A)P,
\end{equation}
one has the exact decomposition
\begin{align}
  f-B_{\rm adm}f
  &=
  (I-P)f
  +
  P\bigl(I-\chi(A)^2\bigr)Pf
  \nonumber\\
  &\quad+
  P\chi(A)\bigl(I-e^{-\tau_{\rm adm}A}\bigr)\chi(A)Pf.
  \label{eq:three-source-error-decomposition}
\end{align}
The three terms are respectively:
\begin{align}
  (I-P)f
  &=
  \text{discarded noncoherent component},
  \\
  P\bigl(I-\chi(A)^2\bigr)Pf
  &=
  \text{spectral-window loss},
  \\
  P\chi(A)\bigl(I-e^{-\tau_{\rm adm}A}\bigr)\chi(A)Pf
  &=
  \text{finite proper-time smoothing}.
\end{align}

\begin{proof}
Insert \(I=P+(I-P)\):
\begin{equation}
  f-B_{\rm adm}f
  =
  (I-P)f
  +
  P f
  -
  P\chi(A)e^{-\tau_{\rm adm}A}\chi(A)P f.
\end{equation}
Since \(P\), \(\chi(A)\), and \(e^{-\tau_{\rm adm}A}\) commute by functional calculus on the
spectral sector, write
\begin{align}
  P f
  -
  P\chi(A)e^{-\tau_{\rm adm}A}\chi(A)P f
  &=
  P f
  -
  P\chi(A)^2P f
  \nonumber\\
  &\quad+
  P\chi(A)^2P f
  -
  P\chi(A)e^{-\tau_{\rm adm}A}\chi(A)P f
  \\
  &=
  P\bigl(I-\chi(A)^2\bigr)P f
  \nonumber\\
  &\quad+
  P\chi(A)\bigl(I-e^{-\tau_{\rm adm}A}\bigr)\chi(A)P f.
\end{align}
This gives \eqref{eq:three-source-error-decomposition}.
\end{proof}

In the common sharp-window case where \(f\in\Ran P\) and \(\chi(A)f=f\), the first two
terms vanish and only finite proper-time smoothing remains. If the relevant spectral support
of \(f\) lies in
\begin{equation}
  \sigma(A|_{\Ran P})\subset[0,\Lambda_{\rm coh}],
\end{equation}
then
\begin{equation}
  \|f-B_{\rm adm}f\|
  \le
  \bigl(1-e^{-\tau_{\rm adm}\Lambda_{\rm coh}}\bigr)\|f\|.
  \label{eq:finite-width-bound}
\end{equation}
For
\begin{equation}
  \tau_{\rm adm}\Lambda_{\rm coh}\ll1,
\end{equation}
this gives
\begin{equation}
  \|f-B_{\rm adm}f\|
  \le
  \tau_{\rm adm}\Lambda_{\rm coh}\|f\|
  +
  O(\tau_{\rm adm}^2\Lambda_{\rm coh}^2)\|f\|.
  \label{eq:small-tau-bound}
\end{equation}

\begin{proposition}[Resolution criterion from finite-width error]
\label{prop:resolution-criterion}
Let \(O\) be an observable whose relevant spectral support is bounded by
\(\Lambda_{\rm obs}\). If the state is already in the accepted coherent window and
\begin{equation}
  \tau_{\rm adm}\Lambda_{\rm obs}\ll \eta,
  \label{eq:resolution-eta}
\end{equation}
where \(\eta\) is the experimental tolerance, then the finite coherent kernel and the sharp
delta idealization are indistinguishable for that observable up to tolerance \(\eta\).
\end{proposition}

\begin{proof}
Under the accepted-window assumptions, apply \eqref{eq:finite-width-bound} on the
observable's retained spectral support and use \(1-e^{-x}\le x\) for \(x\ge0\).
\end{proof}

This is the operational content of pointlikeness in MTT. A particle appears pointlike when
all accessible observables lie far below the coherent cutoff scale and when projection/window
loss is negligible in the sector being probed.

\subsection{Point sources, propagators, and contact idealizations}
\label{subsec:point-sources-propagators}

The same local shadow appears in several familiar places.

A point source,
\begin{equation}
  J(x)=q\delta(x-x_0),
\end{equation}
is replaced by the finite coherent source
\begin{equation}
  J_{\rm adm}(x)
  =
  q\Kadm(x,x_0).
\end{equation}
A point-source Green equation,
\begin{equation}
  LG(x,y)=\delta(x-y),
\end{equation}
is replaced by
\begin{equation}
  LG_{\rm adm}(x,y)=\Kadm(x,y).
\end{equation}
In a flat Euclidean chart with \(A=-\Delta\), the corresponding filtered scalar propagator has
the form
\begin{equation}
  \Delta_{\rm adm}(k)
  =
  \frac{e^{-\ta |k|^2}}{|k|^2+m^2}.
  \label{eq:filtered-propagator-particle-section}
\end{equation}
A local contact vertex is similarly replaced by a finite coherent overlap. Schematically,
\begin{equation}
  \int_X \phi(x)^4\,\dd x
\end{equation}
is replaced by an expression built from products of finite kernels,
\begin{equation}
  \int_X
  \prod_{j=1}^4 \Kadm(x,x_j)
  \,\dd x.
\end{equation}

These examples show that the pointlike particle shadow, the point-source shadow, and the
contact-interaction shadow are manifestations of the same delta-limit structure.

\subsection{Interpretive summary}
\label{subsec:particle-summary}

The particle shadow can now be stated precisely.

A finite coherent excitation gives a finite local response:
\begin{equation}
  \Kadm(x,x_0).
\end{equation}
When its width is below resolution, the response is operationally pointlike. When the coherent
bandwidth is idealized to infinity and the proper-time width to zero, the response becomes
the Dirac delta:
\begin{equation}
  \Kadm(x,y)\to\delta(x-y).
\end{equation}
Thus the point particle is not the fundamental object in MTT. It is the local sharp-limit
shadow of a finite coherent-sector excitation.

\begin{equation}
  \boxed{
  \text{particle-like localization}
  =
  \text{delta-limit shadow of finite coherent support}.
  }
\end{equation}

\section{The wave shadow: spectral phase coherence}
\label{sec:wave-shadow}

The preceding section described the local or particle-like shadow of an admissible coherent
kernel. This section develops the complementary spectral or wave-like shadow. The same
coherent-sector structure that appears locally as a finite source kernel appears spectrally as a
retained modal superposition. Wave-like behavior is the preservation and evolution of relative
phase among these retained modes or branches.

\subsection{Damping operator versus phase generator}
\label{subsec:damping-vs-phase}

Two operators play different roles in the construction.

The operator \(A\) is the fixed-point stabilization or damping operator. It appears in the
proper-time filter
\begin{equation}
  e^{-\ta A},
\end{equation}
and controls coherent support, smoothing, and suppression of discarded modes.

Wave-like behavior, by contrast, requires a phase generator on the retained coherent sector.
We therefore introduce a self-adjoint operator
\begin{equation}
  H_{\rm coh}
\end{equation}
on \(\calH_{\rm coh}=\Ran P\), generating unitary coherent evolution
\begin{equation}
  U_t^{\rm coh}
  =
  e^{-itH_{\rm coh}}.
  \label{eq:coh-unitary}
\end{equation}
In applications \(H_{\rm coh}\) may arise as the effective Hamiltonian, a projected wave
operator, or the phase-transport part of a linearized coherent dynamics. The important point
is that \(H_{\rm coh}\) is conceptually distinct from the heat-time damping operator \(A\).

\begin{assumption}[Coherent phase evolution]
\label{ass:coherent-phase-evolution}
The retained coherent sector \(\calH_{\rm coh}\) carries a self-adjoint phase generator
\(H_{\rm coh}\). The corresponding unitary group \(U_t^{\rm coh}=e^{-itH_{\rm coh}}\)
preserves \(\calH_{\rm coh}\) and transports retained relative phases.
\end{assumption}

\begin{remark}[No conflation of smoothing and phase]
\label{rem:no-conflation}
The heat kernel \(e^{-\ta A}\) determines finite coherent support and admissible smoothing.
The unitary \(e^{-itH_{\rm coh}}\) determines phase evolution. The paper does not identify
these two operations. Their coexistence is what allows a finite coherent excitation to have
both local support and wave-like phase behavior.
\end{remark}

\subsection{Coherent modal expansion}
\label{subsec:coherent-modal-expansion}

Let \(\{\varphi_\alpha\}_{\alpha\in I}\) be an orthonormal modal basis for the retained coherent
sector, chosen so that
\begin{equation}
  H_{\rm coh}\varphi_\alpha
  =
  E_\alpha\varphi_\alpha
\end{equation}
in the simplest diagonal case. A coherent state can be written as
\begin{equation}
  \psi_{\rm coh}
  =
  \sum_{\alpha\in I}
  c_\alpha\varphi_\alpha,
  \qquad
  \sum_{\alpha\in I}|c_\alpha|^2=1.
  \label{eq:coherent-state-expansion}
\end{equation}
Under coherent phase evolution,
\begin{equation}
  \psi_{\rm coh}(t)
  =
  U_t^{\rm coh}\psi_{\rm coh}
  =
  \sum_{\alpha\in I}
  c_\alpha e^{-iE_\alpha t}\varphi_\alpha.
  \label{eq:coherent-phase-expansion}
\end{equation}

The spectral expansion of the admissible kernel,
\begin{equation}
  \Kadm(x,y)
  =
  \sum_{n\in\coh}
  \chi(\lambda_n)^2e^{-\ta\lambda_n}
  \phi_n(x)\phi_n^\ast(y),
\end{equation}
is not itself the time-evolution law. It is the finite coherent support structure. When a
phase generator acts on the retained sector, the retained modal components acquire relative
phases, and those relative phases produce interference.

\begin{definition}[Spectral wave shadow]
\label{def:spectral-wave-shadow}
The spectral wave shadow of a coherent excitation is the representation of \(P\Psi\) as a
retained modal superposition together with its coherent phase evolution under
\(U_t^{\rm coh}\).
\end{definition}

\subsection{Interference from off-diagonal coherence}
\label{subsec:interference-off-diagonal}

Wave-like behavior is most transparently expressed in density-matrix language. Let
\(\rho\) be a density operator on the retained coherent sector. In a branch or modal basis
\(\{|a\rangle\}_{a\in I}\), write
\begin{equation}
  \rho_{ab}
  =
  \langle a|\rho|b\rangle.
\end{equation}
The diagonal entries \(\rho_{aa}\) encode branch weights in that basis. The off-diagonal entries
\(\rho_{ab}\) with \(a\neq b\) encode retained relative phase and coherence.

Let \(E\) be an effect corresponding to an observable in which the branches can interfere.
Then
\begin{equation}
  \operatorname{tr}(\rho E)
  =
  \sum_a \rho_{aa}E_{aa}
  +
  \sum_{a\neq b}\rho_{ab}E_{ba}.
  \label{eq:prob-diagonal-offdiagonal}
\end{equation}
The second term is the interference contribution. It vanishes when the off-diagonal
coherences vanish, and it survives when the branch phases remain coherently related.

\begin{proposition}[Interference criterion]
\label{prop:interference-criterion}
Let \(\rho\) be a density operator on a finite branch subspace and let \(E\) be an effect.
If \(E_{ba}\neq0\) for some \(a\neq b\), then the probability \(\operatorname{tr}(\rho E)\)
depends on the off-diagonal coherence \(\rho_{ab}\). In particular, interference is present
precisely when off-diagonal branch coherence contributes to observable probabilities.
\end{proposition}

\begin{proof}
The identity \eqref{eq:prob-diagonal-offdiagonal} follows from expanding the trace in the
branch basis:
\begin{equation}
  \operatorname{tr}(\rho E)
  =
  \sum_{a,b}\rho_{ab}E_{ba}.
\end{equation}
Terms with \(a=b\) are diagonal contributions, while terms with \(a\neq b\) are off-diagonal
contributions. If \(E_{ba}\neq0\), changing \(\rho_{ab}\) while holding diagonal entries fixed
changes the probability. This is exactly phase-sensitive interference in the chosen context.
\end{proof}

\begin{remark}[Wave-like means phase-sensitive, not spatially diffuse]
\label{rem:wave-like-phase-sensitive}
Wave-like behavior does not mean that the excitation is a classical fluid spread through space.
It means that the effective description retains phase-sensitive off-diagonal structure whose
observable probabilities contain cross terms. Spatial diffraction and interference are common
realizations, but the same logic applies to spin, flavor, path, polarization, and other coherent
branch spaces.
\end{remark}

\subsection{Two-branch superposition}
\label{subsec:two-branch-superposition}

The simplest case is a two-branch coherent state
\begin{equation}
  |\psi\rangle
  =
  \alpha|1\rangle+\beta|2\rangle,
  \qquad
  |\alpha|^2+|\beta|^2=1.
  \label{eq:two-branch-state}
\end{equation}
Its density matrix is
\begin{equation}
  \rho
  =
  |\alpha|^2|1\rangle\langle1|
  +
  |\beta|^2|2\rangle\langle2|
  +
  \alpha\beta^\ast|1\rangle\langle2|
  +
  \alpha^\ast\beta|2\rangle\langle1|.
  \label{eq:two-branch-density}
\end{equation}
The interference-bearing terms are the off-diagonal terms
\begin{equation}
  \rho_{12}=\alpha\beta^\ast,
  \qquad
  \rho_{21}=\alpha^\ast\beta.
\end{equation}

If the two branches acquire a relative phase \(\theta\), then
\begin{equation}
  |\psi(\theta)\rangle
  =
  \alpha|1\rangle+\beta e^{i\theta}|2\rangle.
\end{equation}
For an effect \(E\) with \(E_{21}\neq0\), the probability contains terms of the form
\begin{equation}
  2\operatorname{Re}
  \left[
    \alpha\beta^\ast e^{-i\theta}E_{21}
  \right].
  \label{eq:two-branch-interference-term}
\end{equation}
This is the abstract form of interference fringes.

\begin{corollary}[Loss of off-diagonal coherence removes interference]
\label{cor:loss-coherence-removes-interference}
If the off-diagonal entries \(\rho_{12}\) and \(\rho_{21}\) are set to zero while the diagonal
entries are held fixed, then all phase-dependent two-branch interference terms vanish.
\end{corollary}

\begin{proof}
Set \(\rho_{12}=\rho_{21}=0\) in \eqref{eq:prob-diagonal-offdiagonal}. Only diagonal
terms remain, and these do not depend on the relative phase.
\end{proof}

\subsection{Relation to the admissible kernel}
\label{subsec:wave-relation-kernel}

The wave shadow is not independent of the admissible kernel. The finite kernel selects the
retained coherent sector and weights its spectral content:
\begin{equation}
  \Kadm(x,y)
  =
  \sum_{n\in\coh}
  w_n\phi_n(x)\phi_n^\ast(y),
  \qquad
  w_n=\chi(\lambda_n)^2e^{-\ta\lambda_n}.
  \label{eq:kernel-weights-wave}
\end{equation}
A coherent state supported in this sector has the modal form
\begin{equation}
  \psi_{\rm coh}(x,t)
  =
  \sum_{n\in\coh}
  c_n(t)\phi_n(x).
\end{equation}
When the phase generator is diagonal in the same basis,
\begin{equation}
  c_n(t)=c_n(0)e^{-iE_nt}.
\end{equation}
When \(H_{\rm coh}\) is not diagonal in the \(A\)-basis, coherent phase evolution is still
unitary on \(\Ran P\), and interference is described by the density-matrix criterion above.

Thus \(A\) and \(H_{\rm coh}\) need not be identical. The role of \(A\) is to select and weight
admissible coherent support. The role of \(H_{\rm coh}\) is to transport phase within that
support. Wave-like behavior requires both:
\begin{enumerate}[label=(\roman*)]
  \item a retained coherent sector in which off-diagonal relations can exist;
  \item phase evolution preserving those relations until a measurement context damps or
  selects them.
\end{enumerate}

\subsection{Wave packets and unresolved localization}
\label{subsec:wave-packets}

The spectral wave shadow also includes the familiar wave-packet picture. A wave packet is
a coherent superposition of modes whose phases are arranged to give localized support over
some finite region. In MTT terms, this is not a contradiction between wave and particle
pictures. A wave packet is already a mixed local/spectral representation of the same coherent
sector.

Let
\begin{equation}
  \psi(x)
  =
  \sum_{n\in\coh}
  c_n\phi_n(x)
\end{equation}
be concentrated near \(x_0\). It may be spatially localized, yet still wave-like in the sense that
its evolution and detection probabilities depend on phase relations among the \(c_n\). If its
support is much smaller than detector resolution, it is operationally particle-like; if phase
relations are probed by an interferometer, it is operationally wave-like.

\begin{remark}[No contradiction between localization and interference]
\label{rem:no-contradiction-local-interference}
A localized wave packet can still interfere. Conversely, an extended coherent state can still
produce localized detection events. The apparent contradiction arises only if ``localized'' and
``wave-like'' are treated as mutually exclusive ontologies rather than as different projection
contexts.
\end{remark}

\subsection{Interpretive summary}
\label{subsec:wave-summary}

The wave shadow can now be stated precisely.

A finite coherent excitation has a retained spectral representation:
\begin{equation}
  \psi_{\rm coh}
  =
  \sum_{\alpha}c_\alpha\varphi_\alpha.
\end{equation}
A coherent phase generator transports the retained modes:
\begin{equation}
  \psi_{\rm coh}(t)=e^{-itH_{\rm coh}}\psi_{\rm coh}(0).
\end{equation}
Observable probabilities contain interference exactly when off-diagonal branch or modal
coherences contribute:
\begin{equation}
  \operatorname{tr}(\rho E)
  =
  \sum_a \rho_{aa}E_{aa}
  +
  \sum_{a\neq b}\rho_{ab}E_{ba}.
\end{equation}

Thus the wave is not a separate substance added to a particle. It is the phase-sensitive
spectral shadow of the same coherent-sector excitation whose local shadow can appear
pointlike.

\begin{equation}
  \boxed{
  \text{wave-like interference}
  =
  \text{phase-coherent spectral shadow of finite coherent support}.
  }
\end{equation}

\section{Measurement, survivor basins, and admissible branch damping}
\label{sec:measurement-branch-damping}

The previous two sections described the two shadows of a coherent excitation: local
delta-like localization and spectral phase coherence. Measurement connects them. A
measurement context is a device-dependent disturbance followed by projection and
stabilization. It does not convert a wave into a particle. It selects a downstream survivor
basin and, depending on the device, may damp the off-diagonal coherences that made
interference visible.

This section gives the operator form of that statement.

\subsection{Measurement contexts}
\label{subsec:measurement-contexts}

A measurement is specified not only by an abstract observable but by a physical context.
In MTT, a measurement context \(\Meas\) includes:
\begin{equation}
  \Meas
  =
  (A_{\Meas},P_{\Meas},\chi_{\Meas},\tau_{\Meas},
  \{E_i^{(\Meas)}\}_{i\in I},
  \{\mathcal I_i^{(\Meas)}\}_{i\in I},
  \mathfrak B_{\Meas}).
  \label{eq:measurement-context-data}
\end{equation}
Here:
\begin{enumerate}[label=(\roman*)]
  \item \(A_{\Meas}\) is the detector-sector stabilization or disturbance-damping operator;
  \item \(P_{\Meas}\) is the detector coherent-sector projector;
  \item \(\chi_{\Meas}(A_{\Meas})\) is the detector spectral acceptance window;
  \item \(\tau_{\Meas}\) is the detector-sector proper-time/heat-time scale;
  \item \(\{E_i^{(\Meas)}\}_{i\in I}\) is the finite effect family associated with the possible
  records;
  \item \(\{\mathcal I_i^{(\Meas)}\}_{i\in I}\) is an instrument completion of those effects;
  \item \(\mathfrak B_{\Meas}=\{B_i^{(\Meas)}\}_{i\in I}\) is a proposed corresponding survivor-basin
  partition.
\end{enumerate}

Different measurement devices can have different data. Thus different devices can disturb
the same incoming coherent excitation in different ways. This is not an imperfection added
after the fact. It is part of what defines the measurement context.

\begin{definition}[Closed finite measurement]
\label{def:closed-finite-measurement}
A finite measurement context \(\Meas\) is closed on the retained coherent sector if its effects
satisfy
\begin{equation}
  E_i^{(\Meas)}\ge0,
  \qquad
  \sum_{i\in I}E_i^{(\Meas)}=I_{\rm coh}
  \label{eq:closed-discrete-povm}
\end{equation}
in the discrete-outcome case, or
\begin{equation}
  E_x^{(\Meas)}\ge0,
  \qquad
  \int E_x^{(\Meas)}\,\dd x=I_{\rm coh}
  \label{eq:closed-continuous-povm}
\end{equation}
in the continuous-outcome case.
\end{definition}

For a normalized retained-sector state \(\rho\), the probability of a discrete outcome \(i\) is
\begin{equation}
  p_i^{(\Meas)}
  =
  \tr(\rho E_i^{(\Meas)}),
  \label{eq:discrete-outcome-prob}
\end{equation}
and closedness gives
\begin{equation}
  \sum_i p_i^{(\Meas)}
  =
  \tr(\rho I_{\rm coh})
  =
  1.
\end{equation}

\begin{remark}[Open measurement contexts]
\label{rem:open-measurement-contexts}
If the detector has unobserved leakage channels, the effect family may satisfy
\begin{equation}
  \sum_i E_i^{(\Meas)}\le I_{\rm coh}.
\end{equation}
The missing probability corresponds to outcomes outside the recorded detector context. The
closed case is the appropriate idealization when discussing ordinary normalized measurement
statistics.
\end{remark}

\subsection{Exact records as stabilized survivor basins}
\label{subsec:exact-records}

A detector record is exact relative to its measurement context because the device stabilizes
one branch of the basin partition. This does not mean the device reveals a context-free
pre-existing classical value. It means that, after disturbance and re-coherence, the downstream
state lies in one basin \(B_i^{(\Meas)}\).

Schematically,
\begin{equation}
  \Psi
  \xrightarrow{\text{device disturbance}}
  \Psi'
  \xrightarrow{\text{projection}}
  P_{\Meas}\Psi'
  \xrightarrow{\text{stabilization}}
  B_i^{(\Meas)}.
  \label{eq:measurement-stabilization-chain}
\end{equation}

The stabilized record is discrete even when the pre-measurement coherent excitation was
represented by a continuous wave amplitude. In this sense, particle-like detector clicks are
not primitive point events. They are stable finite records produced by a detector-context
projection.

\begin{definition}[Survivor-basin record]
\label{def:survivor-basin-record}
A survivor-basin record for a measurement context \(\Meas\) is an outcome \(i\in I\) such that
the post-disturbance projected state enters and remains in the basin \(B_i^{(\Meas)}\) under
the local stabilization dynamics.
\end{definition}

\subsection{Outcome probabilities versus coherence visibility}
\label{subsec:outcome-vs-visibility}

The measurement process has two logically distinct parts.

First, it has an outcome-selection part. Operationally, the supplied instrument gives
\(p_i=\tr\mathcal I_i(\rho)=\tr(\rho E_i)\). The basin-measure account proposes the relative
measures
\begin{equation}
  \mathbb P(i)
  =
  \frac{\mu(B_i^{(\Meas)})}
  {\sum_j\mu(B_j^{(\Meas)})},
  \label{eq:basin-probabilities-measurement}
\end{equation}
where \(\mu\) is the admissibility-weighted measure on the relevant ensemble. This equals the
Born probability only when the basin--Born bridge criterion holds.

Second, it has a coherence-visibility part. Before or during final selection, branch coherences
may survive, partially survive, or be suppressed. These off-diagonal terms determine whether
interference is visible.

Thus one must not identify Born probabilities with visibility damping. The former concern
which basin is selected. The latter concerns how much off-diagonal coherence remains between
branches.

\begin{align}
  &\boxed{\text{a valid basin--Born bridge reproduces outcome frequency}},
  \\
  &\boxed{\text{branch damping controls visibility}}.
\end{align}

This distinction is especially important in the double-slit experiment. Each detection event
is localized, but the ensemble distribution can still display interference if branch coherence
survives until detection.

\subsection{Branch density matrices}
\label{subsec:branch-density-matrices}

Let \(\{|a\rangle\}_{a=1}^N\) be a finite branch basis selected by a measurement context.
A retained-sector density matrix has entries
\begin{equation}
  \rho_{ab}
  =
  \langle a|\rho|b\rangle.
\end{equation}
The diagonal terms \(\rho_{aa}\) give branch weights in this basis. The off-diagonal terms
\(\rho_{ab}\), \(a\neq b\), carry branch coherence.

A measurement device that distinguishes branches acts by suppressing the off-diagonal
entries. The simplest branch-damping form is
\begin{equation}
  \rho_{ab}
  \longmapsto
  D_{ab}^{(\Meas)}\rho_{ab}.
  \label{eq:branch-damping-map-entry}
\end{equation}
In matrix notation,
\begin{equation}
  \rho
  \longmapsto
  \mathcal E_D(\rho)
  =
  D\circ \rho,
  \label{eq:schur-map}
\end{equation}
where \(\circ\) denotes entrywise, or Schur, multiplication.

Not every matrix \(D\) gives a physically admissible damping map. The next theorem states
the consistency condition.

\begin{theorem}[Schur-channel consistency]
\label{thm:schur-channel}
Let \(D=(D_{ab})_{a,b=1}^N\) be a positive semidefinite matrix satisfying
\begin{equation}
  D_{aa}=1
  \qquad
  \text{for all }a.
  \label{eq:D-diagonal-one}
\end{equation}
Then
\begin{equation}
  \mathcal E_D(\rho)=D\circ\rho
\end{equation}
is completely positive and trace preserving on \(N\times N\) density matrices.
\end{theorem}

\begin{proof}
Since \(D\succeq0\), there exist vectors \(v_a\) in an auxiliary Hilbert space such that
\begin{equation}
  D_{ab}=\langle v_b,v_a\rangle.
\end{equation}
Define a Stinespring isometry
\begin{equation}
  V|a\rangle=|a\rangle\otimes v_a.
\end{equation}
The condition \(D_{aa}=1\) implies \(\|v_a\|=1\), hence \(V^\ast V=I\). Then
\begin{align}
  \bigl(\operatorname{tr}_{\rm aux}(V\rho V^\ast)\bigr)_{ab}
  &=
  \rho_{ab}\langle v_b,v_a\rangle
  \\
  &=
  D_{ab}\rho_{ab}.
\end{align}
Thus
\begin{equation}
  \operatorname{tr}_{\rm aux}(V\rho V^\ast)=D\circ\rho.
\end{equation}
The map is therefore completely positive. Since \(D_{aa}=1\), diagonal entries are unchanged,
so
\begin{equation}
  \tr(D\circ\rho)=\tr(\rho).
\end{equation}
Hence the map is trace preserving.
\end{proof}

\begin{remark}[Admissibility condition for damping rules]
\label{rem:admissible-damping-rules}
A proposed branch-damping rule is admissible only if it defines such a positive Schur
multiplier. This prevents arbitrary insertion of visibility factors. The damping matrix must
correspond to a valid reduced dynamics or detector-context projection.
\end{remark}

\subsection{Branch-distance damping and Schoenberg positivity}
\label{subsec:schoenberg-damping}

MTT often writes branch damping in the form
\begin{equation}
  D_{ab}^{(\Meas)}
  =
  \exp[-\tau_{\Meas}\Lambda_{ab}^{(\Meas)}],
  \label{eq:branch-distance-damping}
\end{equation}
where \(\Lambda_{ab}^{(\Meas)}\ge0\) measures detector-induced branch separation,
basin-splitting, or distinguishability.

For this to define a valid Schur channel, the matrix \(D\) must be positive semidefinite.
A sufficient and widely useful condition is conditional negative definiteness of
\(\Lambda=(\Lambda_{ab})\).

\begin{definition}[Conditionally negative definite branch separation]
\label{def:cnd}
A real symmetric matrix \(\Lambda=(\Lambda_{ab})\) with \(\Lambda_{aa}=0\) is conditionally
negative definite if
\begin{equation}
  \sum_{a,b}\overline{c_a}c_b\Lambda_{ab}\le0
  \label{eq:cnd-condition}
\end{equation}
for all complex coefficients \(c_a\) satisfying
\begin{equation}
  \sum_a c_a=0.
\end{equation}
\end{definition}

\begin{proposition}[Schoenberg-type admissibility criterion]
\label{prop:schoenberg}
If \(\Lambda\) is conditionally negative definite, then for every \(\tau\ge0\),
\begin{equation}
  D_{ab}=e^{-\tau\Lambda_{ab}}
\end{equation}
is positive semidefinite and satisfies \(D_{aa}=1\). Hence \(D\circ\rho\) is a completely
positive trace-preserving branch-damping channel.
\end{proposition}

\begin{proof}
This is the finite-dimensional Schoenberg theorem: a real symmetric matrix with zero
diagonal is conditionally negative definite if and only if \(e^{-\tau\Lambda}\) is positive
semidefinite for every \(\tau\ge0\). Since \(\Lambda_{aa}=0\), one has \(D_{aa}=1\). The
claim then follows from \cref{thm:schur-channel}.
\end{proof}

A particularly transparent case is squared Hilbert-space distance. If there exist vectors
\(r_a\) in a real Hilbert space such that
\begin{equation}
  \Lambda_{ab}
  =
  \frac12\norm{r_a-r_b}^2,
  \label{eq:squared-distance-Lambda}
\end{equation}
then \(\Lambda\) is conditionally negative definite, and
\begin{equation}
  D_{ab}
  =
  \exp\left[
  -\frac{\tau}{2}\norm{r_a-r_b}^2
  \right]
  \label{eq:gaussian-branch-damping}
\end{equation}
is a valid damping kernel.

\begin{remark}[Device dependence of \(\Lambda_{ab}^{(\Meas)}\)]
\label{rem:device-dependence-lambda}
The same incoming coherent excitation may yield different branch-separation matrices for
different devices. A weak path marker, a strong absorbing detector, a Stern--Gerlach magnet,
and a recombining interferometer do not define the same branch metric. This is why the
measurement method affects whether the wave-like or particle-like shadow is observed.
\end{remark}

\subsection{Visibility as off-diagonal survival}
\label{subsec:visibility-offdiagonal-survival}

For a two-branch experiment, write
\begin{equation}
  \rho
  =
  \begin{pmatrix}
  \rho_{11} & \rho_{12}\\
  \rho_{21} & \rho_{22}
  \end{pmatrix}.
\end{equation}
A valid two-branch damping channel has
\begin{equation}
  D=
  \begin{pmatrix}
  1 & d\\
  d^\ast & 1
  \end{pmatrix},
  \qquad
  |d|\le1.
\end{equation}
Then
\begin{equation}
  \mathcal E_D(\rho)
  =
  \begin{pmatrix}
  \rho_{11} & d\rho_{12}\\
  d^\ast\rho_{21} & \rho_{22}
  \end{pmatrix}.
\end{equation}

If \(d=1\), branch coherence is fully retained. If \(d=0\), the branch state is reduced to the
corresponding incoherent mixture. Intermediate values \(0<|d|<1\) produce partial visibility.

For real nonnegative damping \(d=D_{12}\in[0,1]\), the interference term is simply multiplied
by \(D_{12}\). This is the operator-theoretic basis of the visibility formula used in the
double-slit section.

\subsection{Measurement as projection-duality switch}
\label{subsec:measurement-projection-duality-switch}

The same coherent excitation can therefore be seen under two limiting measurement regimes.

In a coherence-preserving arrangement,
\begin{equation}
  D_{ab}^{(\Meas)}\approx1
  \qquad
  (a\neq b),
\end{equation}
so off-diagonal terms survive. The wave shadow is visible.

In a strongly distinguishing arrangement,
\begin{equation}
  D_{ab}^{(\Meas)}\approx0
  \qquad
  (a\neq b),
\end{equation}
so off-diagonal terms are suppressed. The downstream record is particle-like.

This does not mean the object changed from wave to particle. It means the device changed
which projection shadow survived in the effective description.

\begin{equation}
  \boxed{
  \text{measurement context controls the survival of branch coherence}.
  }
\end{equation}


\section{The double-slit experiment}
\label{sec:double-slit}

The double-slit experiment is the canonical case in which wave-like and particle-like shadows
appear in the same physical setup. A coherent excitation propagates through two available
branches and produces an interference pattern when no branch-distinguishing record is
stabilized. Yet each detection event on the screen is localized. In MTT this is not a paradox.
The interference pattern is the spectral phase shadow of retained branch coherence, while the
localized screen event is the survivor-basin record produced by the detector context.

This section derives the double-slit formula from the branch-damping and finite-effect
framework of \cref{sec:measurement-branch-damping}.

\subsection{Branch decomposition}
\label{subsec:double-slit-branch-decomposition}

Let \(|1\rangle\) and \(|2\rangle\) denote the two path branches associated with the two slits.
A coherent outgoing state after the slits has the form
\begin{equation}
  |\psi\rangle
  =
  \alpha|\psi_1\rangle+\beta|\psi_2\rangle,
  \label{eq:double-slit-state}
\end{equation}
where \(|\psi_1\rangle\) and \(|\psi_2\rangle\) include the propagation amplitudes from the
respective slits to the detection screen. The coefficients \(\alpha,\beta\) include the relative
weights and phases set by the source and slit geometry. We assume normalization
\begin{equation}
  \langle\psi|\psi\rangle=1.
\end{equation}

The corresponding density matrix is
\begin{align}
  \rho
  &=
  |\alpha|^2|\psi_1\rangle\langle\psi_1|
  +
  |\beta|^2|\psi_2\rangle\langle\psi_2|
  \nonumber\\
  &\quad+
  \alpha\beta^\ast|\psi_1\rangle\langle\psi_2|
  +
  \alpha^\ast\beta|\psi_2\rangle\langle\psi_1|.
  \label{eq:double-slit-density}
\end{align}
The first two terms are branch-diagonal. The last two terms are branch off-diagonal and carry
the relative phase required for interference.

Let \(E_x\) be the finite screen effect associated with a detection record at screen coordinate
\(x\). In a sharp idealization one writes \(E_x=|x\rangle\langle x|\). In MTT \(E_x\) is a
finite detector effect, typically of the form
\begin{equation}
  E_x=|k_x\rangle\langle k_x|,
  \qquad
  k_x(y)=K_{\rm det}(y,x),
  \label{eq:screen-effect}
\end{equation}
or a more general positive effect satisfying the POVM normalization condition.

The screen detection density is
\begin{equation}
  I(x)
  =
  \tr(\rho E_x).
  \label{eq:screen-density}
\end{equation}

\subsection{No which-way device: full branch coherence}
\label{subsec:no-which-way}

When no which-way record is stabilized, the two path branches remain in the same coherent
sector until detection. In the branch-damping notation this corresponds to
\begin{equation}
  D_{12}\approx1.
\end{equation}
Substituting \eqref{eq:double-slit-density} into \eqref{eq:screen-density} gives
\begin{align}
  I(x)
  &=
  |\alpha|^2\langle\psi_1|E_x|\psi_1\rangle
  +
  |\beta|^2\langle\psi_2|E_x|\psi_2\rangle
  \nonumber\\
  &\quad+
  2\operatorname{Re}
  \left[
    \alpha\beta^\ast
    \langle\psi_2|E_x|\psi_1\rangle
  \right].
  \label{eq:double-slit-full-interference}
\end{align}
The last term is the interference term.

In the ideal sharp-position notation, this reduces to the familiar expression
\begin{equation}
  I(x)
  =
  |\alpha\psi_1(x)+\beta\psi_2(x)|^2.
  \label{eq:double-slit-sharp}
\end{equation}
The MTT expression \eqref{eq:double-slit-full-interference} is the finite-effect version of
the same formula.

\begin{remark}[Localized events with an interference distribution]
\label{rem:localized-events-interference-distribution}
There is no contradiction between localized screen records and an interference distribution.
Each individual detection is a stabilized record of the screen detector context. The ensemble
distribution displays interference because the off-diagonal branch coherences survive until
the detection effects are applied.
\end{remark}

\subsection{Which-way device: device-dependent branch damping}
\label{subsec:which-way-device}

Now introduce a which-way device \(\Meas_{\rm ww}\). The device couples differently to the
two path branches. In standard notation one often writes
\begin{equation}
  |\psi_1\rangle|D_0\rangle
  +
  |\psi_2\rangle|D_0\rangle
  \longmapsto
  |\psi_1\rangle|D_1\rangle
  +
  |\psi_2\rangle|D_2\rangle.
  \label{eq:which-way-standard}
\end{equation}
The overlap \(\langle D_2|D_1\rangle\) then controls the residual interference.

In MTT this overlap is interpreted as the downstream signature of a device-dependent
branch-separation process. The device defines a branch-distance scale
\begin{equation}
  \Lambda_{12}^{(\rm ww)}
\end{equation}
and a detector proper-time/heat-time scale \(\tau_{\rm ww}\). The associated damping factor is
\begin{equation}
  D_{12}^{(\rm ww)}
  =
  \exp[-\tau_{\rm ww}\Lambda_{12}^{(\rm ww)}],
  \label{eq:which-way-damping}
\end{equation}
provided the resulting \(D\)-matrix defines a valid Schur channel.

The post-device branch density matrix is
\begin{align}
  \rho'
  &=
  |\alpha|^2|\psi_1\rangle\langle\psi_1|
  +
  |\beta|^2|\psi_2\rangle\langle\psi_2|
  \nonumber\\
  &\quad+
  D_{12}^{(\rm ww)}
  \alpha\beta^\ast|\psi_1\rangle\langle\psi_2|
  +
  D_{21}^{(\rm ww)}
  \alpha^\ast\beta|\psi_2\rangle\langle\psi_1|.
  \label{eq:post-which-way-density}
\end{align}
For real symmetric damping,
\begin{equation}
  D_{12}^{(\rm ww)}=D_{21}^{(\rm ww)}\in[0,1].
\end{equation}

The resulting screen density is
\begin{align}
  I_{\rm MTT}(x)
  &=
  |\alpha|^2\langle\psi_1|E_x|\psi_1\rangle
  +
  |\beta|^2\langle\psi_2|E_x|\psi_2\rangle
  \nonumber\\
  &\quad+
  2D_{12}^{(\rm ww)}
  \operatorname{Re}
  \left[
    \alpha\beta^\ast
    \langle\psi_2|E_x|\psi_1\rangle
  \right].
  \label{eq:double-slit-mtt-damped}
\end{align}

Thus the device controls the survival of the wave-like shadow:
\begin{equation}
  D_{12}^{(\rm ww)}\approx1
  \quad\Rightarrow\quad
  \text{interference visible},
\end{equation}
whereas
\begin{equation}
  D_{12}^{(\rm ww)}\approx0
  \quad\Rightarrow\quad
  \text{interference suppressed}.
\end{equation}

\subsection{Weak, intermediate, and strong which-way regimes}
\label{subsec:which-way-regimes}

The MTT reading naturally yields three regimes.

\paragraph{Weak which-way disturbance.}
If the device couples weakly to the path distinction, then
\begin{equation}
  \tau_{\rm ww}\Lambda_{12}^{(\rm ww)}\ll1,
\end{equation}
so
\begin{equation}
  D_{12}^{(\rm ww)}
  =
  1-\tau_{\rm ww}\Lambda_{12}^{(\rm ww)}
  +O((\tau_{\rm ww}\Lambda_{12}^{(\rm ww)})^2).
\end{equation}
The interference pattern survives with only small visibility loss.

\paragraph{Intermediate disturbance.}
If
\begin{equation}
  \tau_{\rm ww}\Lambda_{12}^{(\rm ww)}\sim1,
\end{equation}
then the interference term is partially suppressed. The ensemble distribution interpolates
between the full interference pattern and the incoherent sum.

\paragraph{Strong which-way disturbance.}
If
\begin{equation}
  \tau_{\rm ww}\Lambda_{12}^{(\rm ww)}\gg1,
\end{equation}
then
\begin{equation}
  D_{12}^{(\rm ww)}\approx0,
\end{equation}
and
\begin{equation}
  I_{\rm MTT}(x)
  \approx
  |\alpha|^2\langle\psi_1|E_x|\psi_1\rangle
  +
  |\beta|^2\langle\psi_2|E_x|\psi_2\rangle.
  \label{eq:double-slit-incoherent-sum}
\end{equation}
The wave-like interference shadow is no longer visible. The downstream description is
particle-like in the sense that a definite path-correlated record has been stabilized.

\begin{remark}[Exact outcome within a device context]
\label{rem:exact-outcome-device-context}
A strong which-way device gives an exact outcome relative to the branch partition it defines.
This does not mean that the particle carried a context-independent classical path value
through the apparatus. It means that the device disturbance split the coherent sector into
path-distinguishable survivor basins and stabilized one of them.
\end{remark}

\subsection{Visibility law}
\label{subsec:visibility-law}

In a symmetric two-slit setup, the intensity can be written in the form
\begin{equation}
  I_0(x)
  =
  \bar I(x)\left[1+V_0\cos\varphi(x)\right],
  \label{eq:ideal-visibility-form}
\end{equation}
where \(V_0\) is the ideal visibility and \(\varphi(x)\) is the relative phase. Under MTT branch
damping, the interference term is multiplied by \(D_{12}\), giving
\begin{equation}
  I_{\rm MTT}(x)
  =
  \bar I(x)\left[1+D_{12}V_0\cos\varphi(x)\right].
  \label{eq:mtt-visibility-form}
\end{equation}
Therefore
\begin{equation}
  V_{\rm MTT}
  =
  D_{12}V_0.
  \label{eq:visibility-law}
\end{equation}

\begin{proposition}[Visibility reduction]
\label{prop:visibility-reduction}
In a symmetric two-branch interference experiment with real branch damping \(D_{12}\in[0,1]\),
the measured visibility is
\begin{equation}
  V_{\rm MTT}=D_{12}V_0.
\end{equation}
\end{proposition}

\begin{proof}
The maximum and minimum of \eqref{eq:mtt-visibility-form} are
\begin{equation}
  I_{\max}=\bar I(1+D_{12}V_0),
  \qquad
  I_{\min}=\bar I(1-D_{12}V_0).
\end{equation}
Thus
\begin{equation}
  \frac{I_{\max}-I_{\min}}{I_{\max}+I_{\min}}
  =
  D_{12}V_0.
\end{equation}
\end{proof}

This formula is the quantitative form of the projection-duality crossover. As the device
becomes more path-distinguishing, \(D_{12}\) decreases and the wave-like shadow disappears
from the ensemble distribution.

\subsection{Screen localization and finite detector kernels}
\label{subsec:screen-localization}

The same experiment also displays particle-like localization at the screen. A screen detection
at location \(x\) is represented by a finite effect
\begin{equation}
  E_x=|k_x\rangle\langle k_x|,
  \qquad
  k_x(y)=K_{\rm det}(y,x).
\end{equation}
If \(K_{\rm det}(y,x)\) is narrow relative to the detector resolution, the record is effectively
pointlike.

In the sharp idealization,
\begin{equation}
  K_{\rm det}(y,x)\to\delta(y-x),
\end{equation}
and
\begin{equation}
  E_x\to |x\rangle\langle x|.
\end{equation}
In the proposed MTT encoding, the finite kernel is the retained object. The delta/PVM form
is its singular measurement limit.

Thus the double-slit experiment contains both shadows at once:
\begin{enumerate}[label=(\roman*)]
  \item the ensemble pattern depends on off-diagonal branch coherence and is wave-like;
  \item each detector record is a finite survivor-basin event and is particle-like.
\end{enumerate}

\subsection{Delayed choice and quantum eraser reading}
\label{subsec:delayed-choice-eraser}

The same formalism also explains delayed-choice and quantum eraser variants. What matters
is not whether a particle retroactively chose a path. What matters is which final measurement
context is imposed on the branch-marker system.

If the marker basis preserves path distinguishability, then the branch separation remains
large and
\begin{equation}
  D_{12}\approx0.
\end{equation}
No unconditional interference appears.

If the marker is measured in a recombining basis that groups the two path alternatives into
coherent subensembles, then conditional interference fringes can reappear within those
post-selected subensembles. In MTT language, the final measurement context changes the
survivor-basin partition. It does not rewrite the past; it changes which downstream shadow
is stabilized and which conditional ensemble is being described.

\begin{remark}[No retrocausality required]
\label{rem:no-retrocausality}
Delayed-choice and eraser experiments do not require retroactive changes of history in this
framework. They require only that the final measurement context determine whether the
available branch information is stabilized as distinguishable records or recombined into a
coherent conditional basis.
\end{remark}

\subsection{Summary of the double-slit result}
\label{subsec:double-slit-summary}

The double-slit experiment is not evidence that a quantum object switches between being a
wave and being a particle. It is evidence that two different projection shadows are available.

With branch coherence preserved,
\begin{equation}
  D_{12}\approx1,
\end{equation}
and the wave-like interference term survives. With path-distinguishing disturbance,
\begin{equation}
  D_{12}\approx0,
\end{equation}
and the ensemble becomes an incoherent sum. With finite screen effects, every individual
record is localized.

Thus:
\begin{equation}
  \boxed{
  \text{interference pattern}
  =
  \text{spectral branch-coherence shadow},
  }
\end{equation}
while
\begin{equation}
  \boxed{
  \text{localized screen click}
  =
  \text{finite detector record}.
  }
\end{equation}

Both arise from the same coherent excitation under different projection aspects.
\section{Further experimental modules}
\label{sec:further-experiments}

The double-slit experiment is the cleanest illustration of projection duality, but it is not a
special case in the sense of requiring a unique mechanism. The same architecture appears in
spin measurements, interferometers, delayed-choice arrangements, matter-wave diffraction,
and multi-particle indistinguishability experiments. In each case, MTT separates three
questions:
\begin{enumerate}[label=(\roman*)]
  \item which coherent branches or modes are retained;
  \item how the measurement device disturbs and separates those branches;
  \item which survivor-basin record is stabilized.
\end{enumerate}
This section records several standard modules. They are not new phenomenological fits; they
are consistency checks showing that the same projection-duality structure covers familiar
wave-particle experiments.

\subsection{Stern--Gerlach splitting}
\label{subsec:stern-gerlach}

A Stern--Gerlach apparatus measures spin by coupling spin branches to spatially separated
trajectories. The measurement context is determined by the magnet orientation \(\hat n\).
For spin \(1/2\), the branch basis is
\begin{equation}
  \{|\uparrow_{\hat n}\rangle,|\downarrow_{\hat n}\rangle\}.
\end{equation}
An incoming spin state can be written
\begin{equation}
  |\psi\rangle
  =
  \alpha|\uparrow_{\hat n}\rangle
  +
  \beta|\downarrow_{\hat n}\rangle.
  \label{eq:sg-incoming}
\end{equation}
The Stern--Gerlach field correlates these spin branches with different spatial wave packets:
\begin{equation}
  |\psi\rangle|\Phi_0\rangle
  \longmapsto
  \alpha|\uparrow_{\hat n}\rangle|\Phi_+\rangle
  +
  \beta|\downarrow_{\hat n}\rangle|\Phi_-\rangle.
  \label{eq:sg-branch-correlation}
\end{equation}

In MTT, the magnet plus downstream detector define a measurement context
\begin{equation}
  \Meas_{\rm SG}(\hat n).
\end{equation}
The corresponding branch damping acts on the spin off-diagonal term:
\begin{equation}
  \rho_{\uparrow\downarrow}
  \longmapsto
  D_{\uparrow\downarrow}^{(\hat n)}\rho_{\uparrow\downarrow}.
  \label{eq:sg-damping}
\end{equation}
For a strong Stern--Gerlach measurement,
\begin{equation}
  D_{\uparrow\downarrow}^{(\hat n)}\approx0,
\end{equation}
and the apparatus stabilizes one of two spatially separated survivor basins:
\begin{equation}
  B_+^{(\hat n)}
  \quad\text{or}\quad
  B_-^{(\hat n)}.
\end{equation}

\begin{proposition}[Stern--Gerlach records as context-selected basins]
\label{prop:sg-context-basins}
In a Stern--Gerlach measurement, the exact up/down record is exact relative to the device
orientation \(\hat n\). Changing \(\hat n\) changes the branch basis and hence changes the
survivor-basin partition selected by the device.
\end{proposition}

\begin{proof}
The branch decomposition \eqref{eq:sg-incoming} is defined by the eigenbasis of
\(\sigma\cdot\hat n\). A different orientation \(\hat n'\) defines a different eigenbasis and a
different splitting of the incoming coherent spin state. The measurement device correlates
the selected basis with spatially separated packets and detector basins. Thus the stabilized
record is definite inside the context \(\Meas_{\rm SG}(\hat n)\), not independent of the
chosen orientation.
\end{proof}

This illustrates why measurement outcomes are exact without being context-free. The device
does not merely read a pre-existing classical spin vector. It imposes a branch basis and
stabilizes one branch.

\subsection{Mach--Zehnder interferometer}
\label{subsec:mach-zehnder}

A Mach--Zehnder interferometer is a controlled two-branch experiment. After the first beam
splitter, the state is
\begin{equation}
  |\psi\rangle
  =
  \frac{1}{\sqrt2}
  \left(
  |1\rangle+e^{i\varphi}|2\rangle
  \right),
  \label{eq:mz-state}
\end{equation}
where \(\varphi\) is the relative phase between the two arms.

If the arms are recombined without which-path marking, the output probabilities are
phase-sensitive:
\begin{equation}
  p_\pm
  =
  \frac12(1\pm V_0\cos\varphi),
  \label{eq:mz-ideal}
\end{equation}
with ideal visibility \(V_0\le1\).

A which-path marker or absorber introduces branch separation. In MTT this is represented by
a damping factor
\begin{equation}
  D_{12}^{(\rm MZ)}
  =
  \exp[-\tau_{\rm MZ}\Lambda_{12}^{(\rm MZ)}],
\end{equation}
assuming admissibility of the corresponding Schur channel. The output probabilities become
\begin{equation}
  p_\pm^{\rm MTT}
  =
  \frac12
  \left(
  1\pm D_{12}^{(\rm MZ)}V_0\cos\varphi
  \right).
  \label{eq:mz-mtt}
\end{equation}

The Mach--Zehnder setup therefore makes projection duality tunable. Recombination without
branch marking displays the wave shadow. Strong branch marking suppresses the interference
term and yields path-like records.

\subsection{Delayed choice and quantum eraser}
\label{subsec:delayed-choice-module}

Delayed-choice and quantum eraser experiments are often described as paradoxical because
the final measurement arrangement appears to decide whether the system behaved like a wave
or a particle. In MTT, the issue is not retrospective behavior. The issue is which final
survivor-basin partition is imposed on the joint branch-marker system.

Let the path branches be correlated with marker states:
\begin{equation}
  |\Psi\rangle
  =
  \frac{1}{\sqrt2}
  \left(
  |\psi_1\rangle|M_1\rangle
  +
  e^{i\varphi}|\psi_2\rangle|M_2\rangle
  \right).
  \label{eq:eraser-joint-state}
\end{equation}
If the marker is measured in the distinguishable basis
\begin{equation}
  \{|M_1\rangle,|M_2\rangle\},
\end{equation}
then path information is stabilized and the unconditional interference is suppressed.

If instead the marker is measured in a recombining basis such as
\begin{equation}
  |M_\pm\rangle
  =
  \frac{1}{\sqrt2}
  (|M_1\rangle\pm |M_2\rangle),
  \label{eq:eraser-marker-basis}
\end{equation}
then conditional subensembles can recover phase-sensitive fringes. The final measurement
context has changed the basin partition:
\begin{equation}
  \{B_1,B_2\}
  \quad\longrightarrow\quad
  \{B_+,B_-\}.
\end{equation}

\begin{proposition}[No-retrocausal eraser reading]
\label{prop:no-retrocausal-eraser}
In the MTT reading, delayed-choice and eraser phenomena require no retroactive change of a
past path. They reflect the fact that different final measurement contexts stabilize different
survivor-basin partitions of the joint path-marker state.
\end{proposition}

\begin{proof}
The joint state \eqref{eq:eraser-joint-state} contains both path and marker degrees of
freedom. Measuring the marker in the path-distinguishing basis stabilizes basins correlated
with \(M_1\) and \(M_2\), suppressing path interference in the unconditional ensemble.
Measuring in the recombining basis \eqref{eq:eraser-marker-basis} stabilizes different
conditional basins \(B_\pm\), within which relative phase information can reappear as
conditional fringes. No past event is changed; the final projection context determines the
downstream conditional shadow.
\end{proof}

\subsection{Matter-wave diffraction}
\label{subsec:matter-wave-diffraction}

Electron, neutron, atom, and molecule diffraction experiments show that massive objects
produce interference and diffraction patterns while still being detected as localized events.
This is a direct example of the distinction between ensemble wave shadow and individual
particle shadow.

In a diffraction grating experiment, the grating prepares a coherent superposition of many
path or momentum branches:
\begin{equation}
  |\psi\rangle
  =
  \sum_{a} c_a|a\rangle.
\end{equation}
Free propagation transports the relative phases among the branches. The detection screen
then applies finite effects \(E_x\), giving
\begin{equation}
  p(x)
  =
  \sum_{a,b}
  c_a c_b^\ast
  \langle b|E_x|a\rangle.
  \label{eq:diffraction-prob}
\end{equation}
The off-diagonal terms \(a\neq b\) generate diffraction structure.

If a gas, thermal bath, or monitoring field distinguishes the branches, the density matrix is
modified by an admissible damping matrix:
\begin{equation}
  \rho_{ab}
  \mapsto
  D_{ab}\rho_{ab}.
\end{equation}
The diffraction pattern loses contrast as the off-diagonal terms are suppressed.

The important point is that increasing mass or complexity does not by itself remove
wave-like behavior. What matters is whether the relevant branch coherence remains admissible
relative to environmental and detector disturbances. This is why matter-wave experiments
can display interference for large molecules when isolation is sufficient.

\subsection{Hong--Ou--Mandel-type indistinguishability}
\label{subsec:hom}

The branch-coherence logic is not restricted to single-particle path interference. It also
applies to multi-particle alternatives in configuration space. Hong--Ou--Mandel-type
experiments are governed by indistinguishability between two-particle alternatives.

At a beam splitter, two alternatives may lead to the same detection pattern. If the alternatives
are indistinguishable, their amplitudes interfere. If timing, polarization, frequency, or other
markers distinguish them, the off-diagonal coherence between alternatives is damped.

Let \(a,b\) label the relevant two-particle alternatives. The MTT damping rule has the same
form:
\begin{equation}
  \rho_{ab}
  \mapsto
  D_{ab}^{(\rm HOM)}\rho_{ab},
\end{equation}
where \(D_{ab}^{(\rm HOM)}\) depends on the distinguishability introduced by the device and
input preparation. Perfect indistinguishability corresponds to
\begin{equation}
  D_{ab}^{(\rm HOM)}\approx1,
\end{equation}
while strong distinguishing information gives
\begin{equation}
  D_{ab}^{(\rm HOM)}\approx0.
\end{equation}

Thus particle-like coincidence records and wave-like multi-particle interference again arise
from the same structure: coherent alternatives plus finite measurement effects.

\subsection{Quantum Zeno and anti-Zeno regimes}
\label{subsec:zeno-module}

Repeated measurement provides another illustration of device-dependent disturbance. A
rapid sequence of weak or strong projections can inhibit or accelerate transitions depending
on how the disturbance interacts with the coherent basin geometry.

Let \(P_B\) denote the projector or finite effect associated with remaining in a basin \(B\).
Repeated monitoring over intervals \(\Delta t\) acts schematically as
\begin{equation}
  \rho
  \longmapsto
  \mathcal M_{\Delta t}
  \bigl(
  U_{\Delta t}\rho U_{\Delta t}^\ast
  \bigr),
\end{equation}
where \(\mathcal M_{\Delta t}\) is the finite measurement channel induced by the device.
If repeated disturbance continually projects the state back into the same basin, transitions
are suppressed. This is the Zeno regime. If the disturbance injects fluctuations into a
nearby transition layer or basin boundary, it can enhance switching. This is the anti-Zeno
regime.

In MTT terms, Zeno and anti-Zeno behavior are not separate postulates. They are different
regimes of the same disturbance--damping--stabilization balance.

\subsection{Ramsey and spin-echo reversibility}
\label{subsec:ramsey-echo}

Ramsey interferometry and spin-echo experiments probe how coherences dephase and then
partially rephase under controlled operations. These experiments are important because they
show that not all loss of visible interference is irreversible selection.

In MTT language, dephasing without basin capture corresponds to coherent phase dispersion
inside an admissible sector. Echo pulses can reverse or refocus this dispersion. By contrast,
true survivor-basin stabilization is noninvertible at the effective level: once the measurement
context has stabilized a record and discarded incompatible coherences, the original branch
phase information is not recoverable within the same downstream description.

Thus Ramsey and echo experiments help separate two mechanisms:
\begin{enumerate}[label=(\roman*)]
  \item reversible phase dispersion within the coherent sector;
  \item irreversible basin selection after admissibility-breaking disturbance.
\end{enumerate}
This distinction is important for preventing the MTT measurement account from collapsing
ordinary reversible dephasing and genuine record formation into the same category.

\subsection{Summary of experimental modules}
\label{subsec:experiment-modules-summary}

Across these examples the same structure repeats:
\begin{equation}
  \begin{aligned}
  \text{coherent branches}
  &\longrightarrow \text{device-dependent disturbance}\\
  &\longrightarrow \text{branch damping or basin selection}
  \longrightarrow \text{localized records}.
  \end{aligned}
\end{equation}

The wave-like shadow is visible when branch coherence survives:
\begin{equation}
  D_{ab}\approx1.
\end{equation}
The particle-like shadow dominates when the device distinguishes and stabilizes branches:
\begin{equation}
  D_{ab}\approx0.
\end{equation}
Intermediate regimes yield partial visibility, threshold behavior, Zeno-like inhibition,
anti-Zeno switching, or conditional recovery of fringes, depending on the measurement
context.

Thus the double-slit experiment is not an isolated mystery. It is one instance of a general
projection-duality architecture.


\section{Projection-duality theorem}
\label{sec:projection-duality-theorem}

We now collect the previous constructions into the central theorem of the paper. The theorem
is stated in two layers. The first layer is mathematical: a single admissible coherent operator
has both a local kernel representation and a spectral modal representation, and its sharp
limit recovers the identity kernel. The second layer is physical: under coherent phase
evolution and finite measurement selection, these two representations appear respectively as
particle-like localization and wave-like interference.

\subsection{Mathematical kernel-duality theorem}
\label{subsec:kernel-duality-theorem}

\begin{theorem}[Kernel dual-representation theorem]
\label{thm:kernel-dual-representation}
Let \(X\) be compact and let \(A\ge0\) be a self-adjoint elliptic operator with compact
resolvent. Let \(P\) be the Riesz spectral projector onto an isolated coherent spectral
cluster, let \(\chi(A)\) be an admissible spectral window subordinate to that cluster, and let
\(\ta>0\). Define
\begin{equation}
  \Badm
  =
  P\chi(A)e^{-\ta A}\chi(A)P.
  \label{eq:theorem-badm}
\end{equation}
Then:
\begin{enumerate}[label=(\roman*)]
  \item \(\Badm\) is a bounded positive operator on \(L^2(X)\).
  \item For \(\ta>0\), \(\Badm\) is smoothing on the retained spectral sector and has a smooth
 Schwartz kernel \(\Kadm(x,y)\), under the regularity assumptions of
Assumption~\ref{ass:fixed-point-spectral-data}.
  \item In the spectral representation,
  \begin{equation}
    \Kadm(x,y)
    =
    \sum_{n\in\coh}
    \chi(\lambda_n)^2 e^{-\ta\lambda_n}
    \phi_n(x)\phi_n^\ast(y).
    \label{eq:theorem-spectral-kernel}
  \end{equation}
  \item In the local representation, \(x\mapsto \Kadm(x,x_0)\) is the finite coherent
  response to an ideal source at \(x_0\).
  \item For an exhausting family of projectors and windows
  \((P_N,\chi_N)\) together with \(\tau_N\downarrow0\) satisfying the hypotheses of
  \cref{thm:delta-limit}, the corresponding kernels converge to
  \(\delta(x-y)\) distributionally.
\end{enumerate}
\end{theorem}

\begin{proof}
Since \(A\) is nonnegative self-adjoint, the spectral theorem defines \(e^{-\ta A}\) and
\(\chi(A)\) by functional calculus. Since \(P\) is a Riesz spectral projector of \(A\), it is a
bounded orthogonal spectral projector and commutes with both \(\chi(A)\) and \(e^{-\ta A}\).
Thus \(\Badm\) is bounded. It is positive because it is diagonal in the spectral representation
with nonnegative eigenvalues
\begin{equation}
  p_n\chi(\lambda_n)^2e^{-\ta\lambda_n}\ge0.
\end{equation}

For \(\ta>0\), elliptic heat-kernel regularity implies that \(e^{-\ta A}\) is smoothing. Since
\(P\) and \(\chi(A)\) are spectral multipliers subordinate to the coherent sector, the product
preserves the retained smooth spectral domain and has the asserted smooth kernel under the
regularity assumptions.

The spectral formula follows directly from applying the product
\(P\chi(A)e^{-\ta A}\chi(A)P\) to the eigenbasis \(A\phi_n=\lambda_n\phi_n\). Only retained
modes with \(p_n=1\) contribute, giving \eqref{eq:theorem-spectral-kernel}.

The local response statement follows from the definition of the Schwartz kernel:
\begin{equation}
  (\Badm\delta_{x_0})(x)=\Kadm(x,x_0)
\end{equation}
in the distributional sense. The sharp delta limit is exactly \cref{thm:delta-limit}.
\end{proof}

\begin{remark}[Why the theorem is not merely interpretive]
\label{rem:not-merely-interpretive}
The theorem does not use the words ``wave'' or ``particle.'' It is an operator statement.
It says that a single admissible coherent operator has both a local kernel representation and
a spectral modal representation. The physical terminology enters only after specifying how
experiments read these representations.
\end{remark}

\subsection{Physical projection-duality theorem}
\label{subsec:physical-projection-duality-theorem}

We now add the physical assumptions needed to interpret the two representations.

\begin{assumption}[Physical reading assumptions]
\label{ass:physical-reading}
In addition to the hypotheses of \cref{thm:kernel-dual-representation}, assume:
\begin{enumerate}[label=(\roman*)]
  \item the retained sector carries a self-adjoint coherent phase generator \(H_{\rm coh}\),
  giving unitary phase evolution \(U_t^{\rm coh}=e^{-itH_{\rm coh}}\);
  \item measurement contexts are represented by finite POVMs or effect densities on the
  retained sector;
  \item each operational measurement includes an instrument completing those effects;
  \item branch damping induced by a measurement context is represented by a valid Schur
  channel \(\rho\mapsto D\circ\rho\), with \(D\succeq0\) and \(D_{aa}=1\);
  \item an explicit intertwiner identifies instrument outcomes with survivor basins of the
  measurement context.
\end{enumerate}
\end{assumption}

\begin{theorem}[Conditional wave--particle operator encoding]
\label{thm:wave-particle-projection-duality}
Under Assumptions~\ref{ass:fixed-point-spectral-data},
\ref{ass:coherent-phase-evolution}, and
\ref{ass:physical-reading},
a finite coherent-sector excitation has two downstream shadows:
\begin{enumerate}[label=(\roman*)]
  \item a particle-like local shadow, obtained when the finite local response kernel is probed
  below detector resolution and idealized by the delta limit;
  \item a wave-like spectral shadow, obtained when retained modal or branch phases evolve
  coherently and contribute off-diagonal interference terms to observable probabilities.
\end{enumerate}
Measurement contexts interpolate between these representations by damping or preserving
off-diagonal branch coherences and by applying an instrument. Under the stated basin
intertwiner, instrument outcomes also label survivor-basin records. Thus the supplied data
give one consistent operator encoding of wave-like interference and particle-like localization.
\end{theorem}

\begin{proof}
By \cref{thm:kernel-dual-representation}, the admissible coherent operator has a local kernel
representation and a spectral modal representation.

For the particle-like statement, fix a detector context whose finite effects \(E_x\) have
localization width below the detector resolution. Then the local kernel response
\(\Kadm(x,x_0)\), or the corresponding detector effect \(E_x\), is operationally indistinguishable
from a sharp position effect. By \cref{thm:delta-limit}, the exact sharp position description is
obtained as the distributional delta limit. Hence particle-like localization is the local sharp
shadow of the finite coherent kernel.

For the wave-like statement, Assumption~\ref{ass:coherent-phase-evolution} gives unitary phase transport
on the retained sector. By \cref{prop:interference-criterion}, observable probabilities are
phase-sensitive exactly when off-diagonal branch or modal coherences contribute. Therefore
the spectral representation of the retained coherent sector gives wave-like interference
whenever those coherences survive.

For the measurement statement, \cref{thm:schur-channel} shows that admissible branch
damping is a valid quantum channel when \(D\succeq0\) and \(D_{aa}=1\). If \(D_{ab}\approx1\)
for relevant branches, off-diagonal coherence survives and the wave-like shadow is visible.
If \(D_{ab}\approx0\), off-diagonal coherence is suppressed and the downstream description
is a mixture of branch records. The assumed instrument--basin intertwiner identifies the
operational record with survivor-basin stabilization.

Thus the same declared coherent-sector model represents wave-like or particle-like behavior
according to the measurement context and surviving coherences. This proves the conditional
encoding statement. It does not derive the supplied quantum data from MTT.
\end{proof}

\begin{equation}
  \boxed{
  \text{wave--particle duality}
  \quad\leadsto\quad
  \text{a conditional projection-duality encoding}.
  }
\end{equation}

\subsection{Corollaries}
\label{subsec:projection-duality-corollaries}

\begin{corollary}[Recovery of ordinary wave mechanics]
\label{cor:ordinary-wave-mechanics}
If branch damping is absent or negligible,
\begin{equation}
  D_{ab}\approx1,
\end{equation}
and coherent phase evolution is retained, then the MTT description reduces to ordinary
wave-mechanical interference on the coherent sector.
\end{corollary}

\begin{proof}
When \(D_{ab}\approx1\), the Schur channel leaves off-diagonal coherences approximately
unchanged. Probabilities therefore include the usual interference terms. The coherent
phase generator \(H_{\rm coh}\) supplies the standard unitary phase evolution.
\end{proof}

\begin{corollary}[Recovery of point-particle detection]
\label{cor:point-particle-detection}
If detector kernels are narrow compared with the resolution scale and the sharp effect limit
is taken, then finite MTT detector effects reduce to ordinary pointlike detection effects:
\begin{equation}
  E_x\to |x\rangle\langle x|.
\end{equation}
\end{corollary}

\begin{proof}
By the detector-kernel construction, \(E_x=|k_x\rangle\langle k_x|\). In the sharp detector
limit \(k_x(y)\to\delta(y-x)\), giving \(E_x\to |x\rangle\langle x|\).
\end{proof}

\begin{corollary}[Incoherent mixture from strong branch selection]
\label{cor:incoherent-mixture}
If a measurement context strongly distinguishes branches so that
\begin{equation}
  D_{ab}\approx0
  \qquad
  (a\neq b),
\end{equation}
then interference terms vanish and the downstream ensemble is described by the corresponding
incoherent branch mixture.
\end{corollary}

\begin{proof}
The Schur map sends \(\rho_{ab}\mapsto D_{ab}\rho_{ab}\). If \(D_{ab}\approx0\) for
\(a\neq b\), then all off-diagonal branch terms are suppressed. Observable probabilities
therefore contain only diagonal branch contributions.
\end{proof}

\begin{remark}[Interpretive non-conversion reading]
\label{cor:no-primitive-conversion}
Within the stated assumptions, measurement does not require a primitive conversion of a
wave into a particle. It requires only a change in projection context: branch coherences are
preserved, damped, or selected according to the measurement-device data.
\end{remark}

\subsection{Relation to complementarity}
\label{subsec:relation-complementarity}

The conditional encoding refines, rather than discards, the traditional complementarity
intuition. Complementarity says that wave and particle descriptions are revealed by different
experimental arrangements. The MTT encoding represents those arrangements by different
projection and instrument contexts.

An interference setup preserves or recombines branch coherence:
\begin{equation}
  D_{ab}\approx1.
\end{equation}
A which-way setup distinguishes branches and stabilizes path-correlated records:
\begin{equation}
  D_{ab}\approx0.
\end{equation}
Intermediate setups yield partial visibility.

Because projection is generally non-invertible, no single downstream description captures all
upstream coherent structure. A wave-like description and a particle-like description can both
be accurate in their respective contexts without being jointly reducible to a single classical
picture.

\subsection{Interpretive conclusion}
\label{subsec:projection-duality-interpretive-conclusion}

The theorem gives the formal version of the conditional encoding:
\begin{equation}
  \boxed{
  \text{one finite coherent-sector excitation}
  \quad\Longrightarrow\quad
  \begin{cases}
    \text{local delta-like particle shadow},\\
    \text{spectral phase-coherent wave shadow},\\
    \text{finite measurement-selection shadow}.
  \end{cases}
  }
\end{equation}

The ordinary particle picture is recovered when local finite kernels are idealized by Dirac
deltas. The ordinary wave picture is recovered when retained phases evolve coherently and
off-diagonal terms survive. The ordinary measurement picture is recovered when finite
effects are idealized by sharp projectors and survivor-basin selection has stabilized a record.

Thus the supplied operator data can encode wave--particle duality without a primitive
wave-to-particle conversion. Deriving those data from one MTT source remains open.
\section{Relation to standard quantum mechanics and existing accounts}
\label{sec:relation-existing-accounts}

The projection-duality theorem is not intended to replace the working formalism of quantum
mechanics. It is intended to explain why that formalism has both wave-like and particle-like
faces. This section clarifies how the MTT account relates to ordinary quantum mechanics,
Born probabilities, decoherence, quantum field theory, pilot-wave theory, and extended-object
approaches such as string theory.

\subsection{Recovery of standard quantum mechanics}
\label{subsec:recovery-standard-qm}

Standard quantum mechanics is embedded in the appropriate limiting regimes once its state,
Hamiltonian, effects, and channels are supplied.

When coherent phase evolution is retained, a state evolves as
\begin{equation}
  |\psi(t)\rangle
  =
  e^{-itH_{\rm coh}}|\psi(0)\rangle.
  \label{eq:standard-unitary-recovery}
\end{equation}
When finite detector effects are idealized as sharp projectors,
\begin{equation}
  E_i\to \Pi_i,
  \label{eq:sharp-projector-recovery}
\end{equation}
one obtains the usual projective-measurement description. When branch damping is absent,
\begin{equation}
  D_{ab}=1,
\end{equation}
the off-diagonal terms of the density matrix remain intact and ordinary interference is
recovered. When branch damping is strong,
\begin{equation}
  D_{ab}=0
  \qquad
  (a\neq b),
\end{equation}
the density matrix reduces to the corresponding incoherent mixture in the selected branch
basis.

Thus the conditional construction reproduces the standard cases:
\begin{align}
  D_{ab}=1
  &\quad\Longrightarrow\quad
  \text{coherent wave-like interference},
  \\
  D_{ab}=0
  &\quad\Longrightarrow\quad
  \text{incoherent branch mixture},
  \\
  E_x\to |x\rangle\langle x|
  &\quad\Longrightarrow\quad
  \text{sharp pointlike detection}.
\end{align}

The difference is not in the limiting formulas. The difference is in the order of explanation.
In standard quantum mechanics, wave evolution and pointlike measurement are part of the
formal structure. The MTT program proposes to obtain both as downstream limits of finite
coherent projection; this paper encodes but does not derive that proposal.

\subsection{Born probabilities and basin weights}
\label{subsec:born-basin-weights}

The conditional encoding concerns the coexistence of wave-like interference and
particle-like localization. It does not derive the Born rule. Operationally,
\(\mathbb P(i)=\tr(\rho E_i)\). The MTT corpus proposes to reproduce this pairing with
basin-measure structure.

Let \(\{B_i\}_{i\in I}\) be the proposed survivor-basin partition associated with a
measurement context \(\Meas\), and let \(\mu_\rho\) be a state-dependent
admissibility-weighted basin measure. The candidate basin probability is
\begin{equation}
  \mathbb P(i)
  =
  \frac{\mu_\rho(B_i)}{\sum_j\mu_\rho(B_j)}.
  \label{eq:born-basin-probability-existing}
\end{equation}
It answers the operational question only when it equals \(\tr(\rho E_i)\) for every admitted
state, as required by \cref{thm:basin-born-criterion}:
\begin{equation}
  \text{Which stabilized record occurs?}
\end{equation}

By contrast, branch damping answers the question:
\begin{equation}
  \text{How much interference remains between alternatives?}
\end{equation}
The damping factors \(D_{ab}\) multiply off-diagonal density-matrix entries. They control
visibility, not final outcome frequencies by themselves.

This separation is essential. In the double-slit experiment, individual screen events are
distributed according to the screen effects and basin capture. The ensemble pattern contains
fringes only if branch coherences survive:
\begin{equation}
  \rho_{12}\mapsto D_{12}\rho_{12}.
\end{equation}
Thus:
\begin{equation}
  \boxed{
  \text{supplied Born weights govern selection frequency},
  }
  \qquad
  \boxed{
  \text{branch damping governs visibility}.
  }
\end{equation}

\subsection{Decoherence}
\label{subsec:relation-decoherence}

The MTT branch-damping map resembles decoherence:
\begin{equation}
  \rho_{ab}\mapsto D_{ab}\rho_{ab}.
\end{equation}
This resemblance is not accidental. Environmental entanglement is one physical way of
suppressing off-diagonal coherences in a reduced density matrix. Standard decoherence
therefore captures an important part of the wave-to-classical transition.

The proposed MTT completion is broader and more structured. Decoherence explains how
coherence can become inaccessible after tracing out environmental degrees of freedom. A
full MTT completion would additionally require:
\begin{enumerate}[label=(\roman*)]
  \item a finite coherent kernel determining admissible support;
  \item a detector or environment sector that supplies a valid measurement context;
  \item a positive branch-damping channel;
  \item a survivor-basin stabilization mechanism;
  \item a basin-measure account of outcome frequencies.
\end{enumerate}

In this sense, decoherence is a shadow of one part of the MTT measurement architecture:
coherence suppression. It does not by itself provide the full projection-duality structure,
because it does not identify both pointlike localization and wavelike interference as two
representations of one admissible coherent kernel, nor does it by itself select a unique
stabilized outcome basin.

\subsection{Quantum field theory}
\label{subsec:relation-qft}

Quantum field theory already weakens the classical idea of a particle as a tiny billiard ball.
Particles are excitations of fields, and local detector events arise from interactions between
fields and measuring devices. This is close in spirit to the MTT view that particle-like
behavior is not a fundamental point ontology.

The MTT difference is again structural. QFT commonly uses local fields, point insertions,
distributions, and propagators as part of its formal machinery. In the delta-projection
reading, the pointlike structures
\begin{equation}
  \delta(x-y),
  \qquad
  \phi(x),
  \qquad
  \int \phi(x)^4\,\dd x
\end{equation}
are sharp shadows of finite coherent kernels and overlap structures. A filtered propagator
such as
\begin{equation}
  \Delta_{\rm adm}(k)
  =
  \frac{e^{-\ta |k|^2}}{|k|^2+m^2}
\end{equation}
is not introduced merely as a regulator; in the MTT reading it is the propagator-level shadow
of finite coherent support.

Thus MTT does not deny QFT. It attempts to explain why QFT's local distributional structures
are so effective: they are idealized limits of finite coherent projection.

\subsection{Pilot-wave theory}
\label{subsec:relation-pilot-wave}

Pilot-wave theory gives a clear ontology in which particles have definite positions guided by
a wave. In a double-slit experiment, the guiding wave passes through both slits, while the
particle follows a definite trajectory influenced by that wave. This is a genuine common-source
attempt in the sense that it explains interference and localized detection within one theory.

The MTT account differs in ontology. It does not postulate both a point particle and a guiding
wave as separate real ingredients. Instead, it treats the particle-like and wave-like descriptions
as two downstream shadows of one finite coherent-sector excitation:
\begin{equation}
  \text{particle shadow}
  =
  \text{local delta-limit representation},
\end{equation}
\begin{equation}
  \text{wave shadow}
  =
  \text{spectral phase-coherent representation}.
\end{equation}
Thus pilot-wave theory is dual-ontology at the foundational level, whereas the MTT account
is projection-dual at the effective-description level.

\subsection{String and extended-object approaches}
\label{subsec:relation-string}

String theory and other extended-object approaches also reject fundamental point-particle
ontology. In perturbative string theory, point particles are replaced by one-dimensional
objects whose vibrational modes appear as particle species at scales much larger than the
string length. This shares an important intuition with MTT:
\begin{equation}
  \text{point particle}
  =
  \text{low-resolution limit of a non-pointlike structure}.
\end{equation}

The difference is that string theory's common source is an extended object with vibrational
modes, whereas the MTT common source is an admissible coherent projection kernel. MTT's
construction also places measurement and survivor-basin selection at the center of the
wave-particle account. The pointlike shadow arises from the local delta limit of a finite kernel;
the wave-like shadow arises from retained spectral phase; the measurement record arises from
finite survivor-basin stabilization.

Thus string theory is close in the point-particle replacement sense, but it is not primarily a
projection-duality account of measurement and interference. MTT's claim is not that strings
are wrong, but that the wave-particle issue can be reframed more generally as a statement
about projections of coherent support.

\subsection{Complementarity}
\label{subsec:relation-complementarity-existing}

The traditional complementarity view says that wave and particle descriptions are revealed
by mutually exclusive experimental arrangements. MTT preserves this insight while giving it
a structural mechanism.

Different experimental arrangements implement different projection contexts:
\begin{equation}
  \Meas_{\rm interference}
  \quad\text{preserves or recombines branch coherence},
\end{equation}
whereas
\begin{equation}
  \Meas_{\rm which\text{-}way}
  \quad\text{distinguishes branches and stabilizes records}.
\end{equation}
Because projection is generally non-invertible, no single downstream description captures all
upstream coherent structure. The wave and particle descriptions are therefore complementary
not because nature switches between two substances, but because different projection contexts
expose different shadows of the same coherent-sector excitation.

\subsection{Novelty claim}
\label{subsec:novelty-claim}

The novelty of the present account is not that finite kernels exist, nor that interference comes
from phase, nor that measurement devices can decohere systems. All of these are familiar.

The distinctive claim is the unified selection architecture:
\begin{equation}
  \boxed{
  \text{one finite admissible coherent kernel}
  \quad\Longrightarrow\quad
  \begin{cases}
  \text{local delta-like particle shadow},\\
  \text{spectral phase wave shadow},\\
  \text{finite measurement-selection shadow}.
  \end{cases}
  }
\end{equation}
The kernel is conditionally specified by declared fixed-point data and a certified
damping/admissibility time. A first-principles MTT source theorem is still required. Branch damping is not an arbitrary
visibility factor. It must define a valid quantum channel. Outcome probabilities are not
identified with damping factors. They are basin-measure weights.

This combination is the specific MTT contribution:
\begin{center}
\fbox{\parbox{0.86\linewidth}{\centering
Within the conditional encoding, wave-like and particle-like behavior are two
representations of one declared coherent-sector structure.}}
\end{center}





\section{Compatibility with Prior Measurement Papers}
\label{sec:compatibility-prior-papers}

This section records how the present projection-duality theorem fits with two earlier MTT
developments: the basin-measure account of the Born rule and the disturbance-stabilization
account of measurement. The purpose is to avoid conflating three related but distinct
structures:
\begin{enumerate}[label=(\roman*)]
  \item basin-measure probabilities;
  \item branch-coherence visibility;
  \item localized survivor-basin records.
\end{enumerate}

\subsection{Compatibility with the basin-measure Born-rule bridge}
\label{subsec:compatibility-born-bridge}

The basin-measure bridge proposes quantum probabilities as shadows of basin measures on an
admissible coherent sector. In its simplest form, if a measurement context has survivor basins
\begin{equation}
  \{B_i\}_{i\in I},
\end{equation}
and if \(\mu_\rho\) is a state-dependent admissibility-weighted basin measure, the proposed
outcome probability is
\begin{equation}
  \mathbb P(i)
  =
  \frac{\mu(B_i)}{\sum_j\mu(B_j)}.
  \label{eq:compat-born-basin}
\end{equation}
This answers the outcome-frequency question only if it agrees with the operational effect
probability \(\tr(\rho E_i)\).

\begin{theorem}[Exact basin--Born bridge criterion]
\label{thm:basin-born-criterion}
Let \(\{E_i\}\) be a POVM and \(\{B_i\}\) a measurable basin partition. The normalized basin
rule reproduces the Born rule for every admitted state \(\rho\) if and only if there is a
positive normalization \(c(\rho)\) such that
\[
  \mu_\rho(B_i)=c(\rho)\tr(\rho E_i)
  \qquad\text{for every }i.
\]
When \(\mu_\rho(\Omega)=1\), one has \(c(\rho)=1\).
\end{theorem}

\begin{proof}
If the displayed condition holds, normalization cancels \(c(\rho)\) because
\(\sum_i\tr(\rho E_i)=1\). Conversely, equality of the normalized basin and Born
probabilities implies the displayed condition with
\(c(\rho)=\sum_j\mu_\rho(B_j)\). \(\square\)
\end{proof}

The theorem isolates rather than solves the source problem: MTT must construct
\(\mu_\rho\) and the basin/effect correspondence from the same upstream data.
The operational question is:
\begin{equation}
  \text{Which stabilized record occurs, and how often?}
\end{equation}

The present paper answers a different question:
\begin{equation}
  \text{Why can the same excitation produce both interference and localized records?}
\end{equation}
The answer is that the excitation has both a spectral phase shadow and a local delta-like
shadow, and measurement contexts determine which branch coherences survive.

Thus the proposed Born bridge and the projection-duality encoding are complementary. The
former still owes the criterion in \cref{thm:basin-born-criterion}; the latter supplies only
the representation structure explaining how coherent interference can coexist with localized
detection once standard quantum data are supplied.

\begin{proposition}[Separation of probability and visibility]
\label{prop:probability-visibility-separation}
Let a measurement context \(\Meas\) have survivor basins \(B_i\) with basin probabilities
\(\mathbb P(i)\) given by \eqref{eq:compat-born-basin}. Let the same context induce a valid
branch-damping channel
\begin{equation}
  \rho_{ab}\mapsto D_{ab}^{(\Meas)}\rho_{ab}.
\end{equation}
Then \(\mathbb P(i)\) and \(D_{ab}^{(\Meas)}\) answer different structural questions:
\begin{enumerate}[label=(\roman*)]
  \item \(\mathbb P(i)\) determines the frequency of stabilized records;
  \item \(D_{ab}^{(\Meas)}\) determines the survival of interference between branches \(a,b\).
\end{enumerate}
Therefore branch damping must not be identified with the Born rule.
\end{proposition}

\begin{proof}
The basin probability \(\mathbb P(i)\) is defined on the partition of stabilized outcome basins.
It assigns weights to final records. The damping factor \(D_{ab}^{(\Meas)}\) acts on
off-diagonal density-matrix entries before or during stabilization and determines whether
cross terms contribute to observable probabilities. Since diagonal basin weights and
off-diagonal coherence survival act on different components of the effective description,
they are not the same object.
\end{proof}

\subsection{Minimal basin-capture assumptions}
\label{subsec:minimal-basin-capture}

The basin-measure formula requires assumptions. In the setting of this paper, the minimal
ones are:
\begin{enumerate}[label=(\roman*)]
  \item a stationary or experimentally reproducible input ensemble \((\Omega,\mu)\);
  \item a measurable survivor-basin partition
  \begin{equation}
    \Omega=\bigcup_i B_i,
    \qquad
    B_i\cap B_j=\varnothing\quad(i\neq j),
  \end{equation}
  up to measure-zero boundaries;
  \item almost-sure finite capture into exactly one basin under the measurement-context
  stabilization dynamics.
\end{enumerate}
Under these assumptions,
\begin{equation}
  \mathbb P(i)=\frac{\mu(B_i)}{\mu(\Omega)}
\end{equation}
is well-defined. If the partition is not measurable, if capture fails, or if trajectories hover
near basin boundaries for experimentally relevant times, then the simple basin-probability
formula must be replaced by a metastable or protocol-dependent description.

\begin{remark}[Why this matters for duality]
\label{rem:basin-capture-duality}
The double-slit interference pattern concerns ensemble density on a screen. The individual
screen click concerns basin capture by the detector. These are compatible only if one keeps
the ensemble interference structure distinct from the single-record stabilization process.
\end{remark}

\subsection{Compatibility with measurement as disturbance and stabilization}
\label{subsec:compatibility-measurement-disturbance}

The disturbance-stabilization measurement account states that a device perturbs a coherent
mode away from a quiet or balanced submanifold. The system then undergoes local smoothing,
global projection, and stabilization into an admissible branch. In schematic form:
\begin{equation}
  \text{disturbance}
  \longrightarrow
  \text{damping/smoothing}
  \longrightarrow
  \text{projection}
  \longrightarrow
  \text{stabilized outcome}.
  \label{eq:disturbance-stabilization-chain}
\end{equation}

The present paper supplies an operator-theoretic version of this chain for wave-particle
duality:
\begin{equation}
  \rho
  \longmapsto
  D^{(\Meas)}\circ\rho
  \longmapsto
  \{E_i^{(\Meas)}\}
  \longmapsto
  B_i^{(\Meas)}.
  \label{eq:operator-disturbance-chain}
\end{equation}
Here \(D^{(\Meas)}\) records the survival of branch coherence, the effect family
\(\{E_i^{(\Meas)}\}\) records the finite measurement readout, and \(B_i^{(\Meas)}\) records the
survivor basin.

This is why different devices matter. A weak which-way marker, a destructive absorber, a
Stern--Gerlach magnet, and a phase-recombining interferometer impose different disturbance
maps and different survivor-basin partitions.

\subsection{Disturbance strength and visibility knees}
\label{subsec:disturbance-strength-knees}

Let \(\delta\ge0\) denote a scalar measure of measurement strength, such as marker coupling,
environmental distinguishability, detector gain, or branch-separation amplitude. A simple
MTT-compatible visibility model is
\begin{equation}
  D_{12}(\delta)
  =
  \exp[-\tau_{\rm det}\Lambda_{12}^{\rm det}(\delta)].
  \label{eq:visibility-delta-model}
\end{equation}
If
\begin{equation}
  \Lambda_{12}^{\rm det}(\delta)
\end{equation}
is small for weak disturbance and grows across a branch-separation threshold, then the
visibility
\begin{equation}
  V(\delta)
  =
  D_{12}(\delta)V_0
\end{equation}
falls smoothly at first and then rapidly near the threshold. This produces a knee-like
transition between wave-like and particle-like behavior.

\begin{remark}[Weak, strong, and threshold regimes]
\label{rem:weak-strong-threshold}
The disturbance-stabilization paper emphasizes that measurement need not be all-or-nothing.
The present formula makes that continuous structure explicit. Weak disturbance leaves
\(D_{12}\approx1\), strong disturbance gives \(D_{12}\approx0\), and threshold regimes can
produce rapid visibility loss, switching, or hysteresis depending on basin geometry.
\end{remark}

\subsection{OU floors and finite detector width}
\label{subsec:ou-floors-detector-width}

The disturbance-stabilization account also interprets uncertainty floors as finite balances
between disturbance and damping. In the present notation, this appears as the finite width of
the detector or coherent kernel. Even after damping, the kernel need not collapse to a delta.
Instead, it may approach a finite Ornstein--Uhlenbeck-like floor set by the balance between
disturbance injection and stabilization.

This supports the hierarchy used throughout the paper:
\begin{equation}
  \text{finite POVM/effect}
  \quad\text{is native},
\end{equation}
while
\begin{equation}
  \text{PVM/delta measurement}
  \quad\text{is a singular idealization}.
\end{equation}

In double-slit language, this means that the localized screen record is never literally a
mathematical point. It is a finite stabilized detector event whose width is below the relevant
resolution scale.

\subsection{What this paper adds}
\label{subsec:what-this-paper-adds}

The earlier measurement papers explain how outcomes stabilize and how probabilities arise
from basins. The present paper adds a specific structural synthesis:
\begin{center}
\emph{The same finite coherent kernel has both a local delta-like representation and a
spectral phase-coherent representation.}
\end{center}
It also adds a channel-consistency condition for interference damping:
\begin{equation}
  D\succeq0,
  \qquad
  D_{aa}=1.
\end{equation}
This ensures that the visibility-reduction map is a valid quantum operation, not merely an
interpretive multiplier.

Thus the paper should be read as a bridge between three parts of MTT:
\begin{enumerate}[label=(\roman*)]
  \item the delta-projection sequence, where pointlike objects are sharp limits of finite kernels;
  \item the measurement sequence, where outcomes are disturbance-stabilized survivor basins;
  \item the basin-measure sequence, where Born probabilities are basin weights.
\end{enumerate}

\subsection{Compatibility summary}
\label{subsec:compatibility-summary}

The relationship can be summarized as:
\begin{align}
  \text{finite coherent kernel}
  &\longrightarrow
  \text{local particle shadow and spectral wave shadow},
  \\
  \text{branch damping}
  &\longrightarrow
  \text{visibility suppression},
  \\
  \text{survivor-basin capture}
  &\longrightarrow
  \text{individual outcome records},
  \\
  \text{basin measure}
  &\longrightarrow
  \text{Born frequencies}.
\end{align}

These are not competing explanations. They are different layers of one MTT measurement
architecture.
\section{Phenomenological consequences and testable scalings}
\label{sec:phenomenological-consequences}

The projection-duality theorem is primarily structural. It identifies wave-like and particle-like
behavior as two shadows of a finite coherent-sector excitation. Nevertheless, the framework
has phenomenological content. Once a sector supplies the relevant operator data, detector
kernel, damping scale, and branch-separation metric, the theory gives definite finite-width and
visibility corrections.

This section records the leading scaling relations. They should not be read as universal
numerical predictions. They are templates: in each physical sector the quantities
\(\ta\), \(\tau_{\Meas}\), \(\Lambda_{\rm eff}\), \(D_{ab}\), and \(\Lambda_{ab}^{(\Meas)}\) must be derived or
bounded from the corresponding coherent-sector and detector-sector data.

\subsection{Pointlikeness bounds}
\label{subsec:pointlikeness-bounds}

The particle-like approximation is valid when the coherent width is below experimental
resolution:
\begin{equation}
  \ell_{\rm coh}\ll \ell_{\rm res}.
  \label{eq:phen-pointlike-length}
\end{equation}
In heat-kernel models,
\begin{equation}
  \ell_{\rm coh}\sim \sqrt{\ta}.
\end{equation}
Equivalently, in energy or momentum terms,
\begin{equation}
  E_{\rm probe}\ll \Lambda_{\rm eff},
  \qquad
  \Lambda_{\rm eff}\sim\ta^{-1/2}.
  \label{eq:phen-pointlike-energy}
\end{equation}

If an experiment observes no deviation from pointlike behavior up to characteristic scale
\(E_{\rm max}\) with tolerance \(\eta\), the generic scaling bound is
\begin{equation}
  \ta E_{\rm max}^2\lesssim \eta.
  \label{eq:phen-tau-bound}
\end{equation}
Thus
\begin{equation}
  \Lambda_{\rm eff}
  \gtrsim
  \frac{E_{\rm max}}{\sqrt{\eta}}.
  \label{eq:phen-lambda-bound}
\end{equation}

This is the same logic used in effective form-factor constraints. The MTT difference is that
\(\ta\) is not introduced merely as a fitted form-factor scale. It is, in principle, fixed by
discarded-sector damping and admissibility:
\begin{equation}
  \ta
  =
  \lambda_\ast^{-1}\log(C_Q/\epsadm).
  \label{eq:phen-tau-fixed}
\end{equation}

\begin{remark}[Interpretation of null results]
\label{rem:null-results-pointlike}
A null result in a high-resolution scattering or localization experiment does not prove that
particles are literal points. In the MTT reading it implies that the finite coherent width is
below the resolution scale tested by the experiment.
\end{remark}

\subsection{Finite propagator corrections}
\label{subsec:finite-propagator-corrections}

In a flat Euclidean coherent chart with \(A=-\Delta\), the finite coherent kernel gives the
filtered scalar propagator
\begin{equation}
  \Delta_{\rm adm}(k)
  =
  \frac{e^{-\ta |k|^2}}{|k|^2+m^2}.
  \label{eq:phen-propagator}
\end{equation}
At low momentum,
\begin{equation}
  e^{-\ta |k|^2}
  =
  1-\ta |k|^2+O(\ta^2|k|^4),
\end{equation}
so
\begin{equation}
  \Delta_{\rm adm}(k)
  =
  \frac{1}{|k|^2+m^2}
  -
  \ta\frac{|k|^2}{|k|^2+m^2}
  +
  O\!\left(
  \frac{\ta^2|k|^4}{|k|^2+m^2}
  \right).
  \label{eq:phen-propagator-expansion}
\end{equation}

The leading correction is controlled by
\begin{equation}
  \ta |k|^2.
\end{equation}
Therefore short-distance deviations from pointlike propagation are expected only when
\begin{equation}
  |k|^2\sim \ta^{-1}.
\end{equation}

\begin{remark}[Gauge and Lorentzian caution]
\label{rem:gauge-lorentzian-caution-phen}
The scalar Euclidean propagator formula is a clean model of finite coherent support. Gauge
theories and Lorentzian scattering require compatibility with gauge identities, positivity,
causality, and unitarity. The projection-duality theorem does not by itself prove those
sector-specific consistency conditions.
\end{remark}

\subsection{Visibility loss}
\label{subsec:visibility-loss-phen}

In a two-branch interference experiment, the MTT visibility law is
\begin{equation}
  V_{\rm MTT}
  =
  D_{12}V_0.
  \label{eq:phen-visibility}
\end{equation}
If the damping factor is
\begin{equation}
  D_{12}
  =
  \exp[-\tau_{\Meas}\Lambda_{12}^{(\Meas)}],
  \label{eq:phen-D}
\end{equation}
then for weak disturbance,
\begin{equation}
  1-D_{12}
  =
  \tau_{\Meas}\Lambda_{12}^{(\Meas)}
  +
  O\!\left((\tau_{\Meas}\Lambda_{12}^{(\Meas)})^2\right).
  \label{eq:phen-visibility-small}
\end{equation}

If an experiment observes no visibility loss above tolerance \(\eta_V\), then
\begin{equation}
  \tau_{\Meas}\Lambda_{12}^{(\Meas)}
  \lesssim
  \eta_V.
  \label{eq:phen-visibility-bound}
\end{equation}
Thus interference visibility bounds the product of measurement-sector time scale and
branch-separation scale.

\begin{remark}[What visibility constraints measure]
\label{rem:visibility-constraints-measure}
Visibility constraints do not directly measure the Born probabilities of final records. They
measure the survival of off-diagonal coherence between alternatives. This is why a high
visibility experiment constrains \(D_{12}\), whereas repeated record frequencies constrain
basin measures.
\end{remark}

\subsection{Disturbance-dependent knees}
\label{subsec:disturbance-knees-phen}

Let \(\delta\) denote a controllable disturbance strength. Examples include which-way marker
coupling, gas pressure in matter-wave interferometry, detector gain, magnetic-field gradient,
or timing distinguishability. Suppose
\begin{equation}
  D_{12}(\delta)
  =
  \exp[-\tau_{\Meas}\Lambda_{12}^{(\Meas)}(\delta)].
  \label{eq:phen-D-delta}
\end{equation}
If \(\Lambda_{12}^{(\Meas)}(\delta)\) grows slowly at first and then rapidly near an admissibility
threshold, the measured visibility
\begin{equation}
  V(\delta)
  =
  V_0D_{12}(\delta)
  \label{eq:phen-V-delta}
\end{equation}
will display a knee-like transition.

This gives a qualitative experimental signature:
\begin{equation}
  \text{smooth weak-disturbance visibility loss}
  \quad\longrightarrow\quad
  \text{rapid threshold suppression}.
\end{equation}
Such knees are expected when the measurement disturbance pushes the state across a
survivor-basin boundary rather than merely adding small reversible phase noise.

\subsection{Reversibility versus basin capture}
\label{subsec:reversibility-phen}

The distinction between reversible dephasing and irreversible basin capture is experimentally
important.

If a disturbance only disperses phases inside the coherent sector, then echo or reversal
protocols may recover visibility:
\begin{equation}
  D_{12}^{\rm effective}\to1
\end{equation}
after rephasing.

If a disturbance stabilizes a survivor-basin record, the lost branch coherence is not recoverable
inside the same downstream description:
\begin{equation}
  D_{12}\approx0
  \quad\text{after record stabilization}.
\end{equation}

Thus MTT predicts a qualitative distinction between:
\begin{enumerate}[label=(\roman*)]
  \item reversible coherence dispersion;
  \item irreversible measurement-induced selection.
\end{enumerate}
Ramsey and echo experiments probe the first; strong which-way measurement probes the
second.

\subsection{Device-dependence as a prediction}
\label{subsec:device-dependence-phen}

Because different devices define different measurement contexts,
\begin{equation}
  \Meas_1\neq \Meas_2,
\end{equation}
they may induce different damping matrices even when aimed at the same nominal observable:
\begin{equation}
  D_{ab}^{(\Meas_1)}
  \neq
  D_{ab}^{(\Meas_2)}.
\end{equation}
Thus MTT predicts that visibility loss is not determined solely by the abstract observable
label. It depends on the device's disturbance channel, detector kernel, and branch-separation
metric.

This is not surprising operationally, but MTT makes it structural. The device is part of the
projection context, not an external passive reader.

\subsection{Summary of testable scalings}
\label{subsec:testable-scalings-summary}

The main phenomenological scalings are:
\begin{align}
  \ell_{\rm coh}
  &\sim
  \sqrt{\ta},
  \\
  \Lambda_{\rm eff}
  &\sim
  \ta^{-1/2},
  \\
  \ta E_{\rm max}^2
  &\lesssim
  \eta
  \quad
  \text{for unresolved pointlike behavior},
  \\
  V_{\rm MTT}
  &=
  D_{12}V_0,
  \\
  D_{12}
  &=
  \exp[-\tau_{\Meas}\Lambda_{12}^{(\Meas)}],
  \\
  \tau_{\Meas}\Lambda_{12}^{(\Meas)}
  &\lesssim
  \eta_V
  \quad
  \text{when visibility loss is unobserved}.
\end{align}

These formulas do not by themselves fix the numerical scales. They become predictions only
when the physical sector supplies the fixed-point data, detector kernel, damping gap, basin
metric, and admissibility tolerance. Their importance is that they identify where measurable
departures from ideal wave or particle shadows should enter.
\section{Failure modes and domain of validity}
\label{sec:failure-modes}

The projection-duality theorem is conditional. It does not assert that every physical regime
admits a clean coherent-sector description, nor that every proposed damping factor or detector
kernel is admissible. This section records the main ways in which the construction can fail.
Making these failure modes explicit is important: it prevents the theorem from becoming a
catch-all interpretation and clarifies what must be checked in any concrete physical sector.

\subsection{Failure of coherent spectral isolation}
\label{subsec:failure-spectral-isolation}

The construction assumes an isolated coherent spectral cluster and a corresponding Riesz
projector
\begin{equation}
  P
  =
  \frac{1}{2\pi i}
  \oint_\Gamma (z-A)^{-1}\,\dd z.
\end{equation}
If the coherent cluster is not isolated, then the contour \(\Gamma\) cannot be chosen without
enclosing discarded spectrum. In that case the retained coherent sector is not sharply defined
by the spectral data of \(A\).

This can occur near gap closure, phase transitions, strong coupling, uncontrolled turbulence,
horizon-like breakdowns, or measurement thresholds. In such regimes, the effective description
may lose stable representability. The failure is not a small correction to the wave-particle
duality theorem; it is a failure of the theorem's hypotheses.

\begin{equation}
  \lambda_\ast\downarrow0
  \quad\Longrightarrow\quad
  \text{coherent-sector projection may cease to be stable}.
\end{equation}

\subsection{Divergent or undefined admissible proper time}
\label{subsec:failure-tau}

The time certified by the damping bound is
\begin{equation}
  \ta
  =
  \lambda_\ast^{-1}\log\frac{C_Q}{\epsadm}.
\end{equation}
This formula requires
\begin{equation}
  \lambda_\ast>0,
  \qquad
  0<\epsadm<1,
  \qquad
  C_Q\ge1.
\end{equation}
If \(\lambda_\ast=0\), the discarded sector is not exponentially damped away from the coherent
sector, and the required admissible time may diverge. If no meaningful basin tolerance
\(\epsadm\) exists, then \(\ta\) is not determined by the damping criterion.

Thus the finite kernel is not guaranteed in every regime. It is selected only when the sector
has a stable damping/admissibility structure.

\subsection{Failure of projector regularity}
\label{subsec:failure-projector-regularity}

The smooth kernel statements require more than a bounded projector. They require \(P\) to
be a spectral/Riesz projector of \(A\), or otherwise to preserve the relevant smooth domain.
If \(P\) is an arbitrary bounded projection that does not respect the spectral regularity of
\(A\), then
\begin{equation}
  P\chi(A)e^{-\ta A}\chi(A)P
\end{equation}
may not have the smooth kernel properties assumed in the local and spectral analysis.

This is why the theorem uses spectral projectors rather than arbitrary coherent labels.

\subsection{Failure of coherent phase evolution}
\label{subsec:failure-phase-evolution}

Wave-like behavior requires a coherent phase generator
\begin{equation}
  H_{\rm coh}
\end{equation}
on the retained sector. If the retained sector does not support approximately unitary phase
transport, then modal coherence may not persist long enough to generate interference.

Such failure may occur when:
\begin{enumerate}[label=(\roman*)]
  \item the retained sector is strongly open;
  \item phase coherence is rapidly lost to uncontrolled discarded modes;
  \item the effective Hamiltonian is not self-adjoint on the retained sector;
  \item the experiment probes beyond the coherent-sector validity regime.
\end{enumerate}

In such cases the particle-like local shadow may remain meaningful while the wave-like
interference shadow is suppressed or absent.

\subsection{Invalid branch-damping matrices}
\label{subsec:failure-invalid-D}

A proposed branch-damping rule
\begin{equation}
  \rho_{ab}\mapsto D_{ab}\rho_{ab}
\end{equation}
is physically admissible only if the Schur map is completely positive and trace preserving.
For finite branch bases this requires
\begin{equation}
  D\succeq0,
  \qquad
  D_{aa}=1.
\end{equation}
If these conditions fail, then the map may send valid density matrices to non-positive
operators. Such a rule cannot represent a physical measurement or disturbance channel.

Similarly, a proposed exponential form
\begin{equation}
  D_{ab}=e^{-\tau\Lambda_{ab}}
\end{equation}
is automatically safe for all \(\tau\ge0\) only under conditions such as conditional negative
definiteness of \(\Lambda\). Otherwise positivity must be checked directly.

\subsection{Unnormalized or leaky detector effects}
\label{subsec:failure-povm}

The finite detector effects must satisfy
\begin{equation}
  \sum_i E_i=I_{\rm coh}
\end{equation}
or
\begin{equation}
  \int E_x\,\dd x=I_{\rm coh}
\end{equation}
for a closed probability-conserving detector on the retained sector. If instead
\begin{equation}
  \sum_i E_i<I_{\rm coh},
\end{equation}
then the detector is open: probability leaks into unrecorded channels. This is not necessarily
a contradiction, but the interpretation changes. One must then include the missing channels
or explicitly work with conditional probabilities.

Thus finite detector kernels do not automatically give normalized measurements. Their POVM
normalization is part of the measurement-context data.

\subsection{Over-resolved coherent width}
\label{subsec:failure-over-resolved}

The particle-like approximation is resolution-dependent. If the experimental probe reaches
the coherent scale
\begin{equation}
  E_{\rm probe}\sim\Lambda_{\rm eff},
\end{equation}
or equivalently
\begin{equation}
  \ell_{\rm res}\sim\ell_{\rm coh},
\end{equation}
then the delta approximation can fail. In that regime the finite kernel should produce
observable deviations from pointlike behavior, such as:
\begin{enumerate}[label=(\roman*)]
  \item softened short-distance response;
  \item finite-width detector effects;
  \item modified contact interactions;
  \item reduced ultraviolet sensitivity;
  \item departure from sharp localization assumptions.
\end{enumerate}

This is not a failure of MTT. It is precisely where MTT predicts that the point-particle shadow
should no longer be treated as exact.

\subsection{Device-context dependence}
\label{subsec:failure-device-context}

Different measurement devices define different contexts:
\begin{equation}
  \Meas_1\neq \Meas_2.
\end{equation}
They may have different branch bases, detector effects, damping matrices, and survivor-basin
partitions. Therefore one cannot compare outcomes from different devices as if they were
readings of the same context-free classical variable.

This is not a loophole. It is part of the theory's content. But in applications it requires care:
the measurement context must be specified before asking which branch basis is being damped
or which basin partition is being sampled.

\subsection{Failure of basin capture}
\label{subsec:failure-basin-capture}

The basin-measure account of outcome probabilities assumes that post-disturbance trajectories
are captured by survivor basins in finite effective time, at least almost surely relative to the
relevant ensemble measure. If trajectories hover near basin boundaries, undergo chaotic
switching, or fail to stabilize, then ordinary discrete outcome statistics may not apply without
additional coarse-graining.

In such cases one may observe:
\begin{enumerate}[label=(\roman*)]
  \item metastability;
  \item threshold knees;
  \item switching noise;
  \item hysteresis;
  \item Zeno or anti-Zeno behavior;
  \item protocol dependence.
\end{enumerate}
These are not outside MTT, but they are outside the simplest stable-basin theorem used in
this paper.

\subsection{Domain of validity}
\label{subsec:domain-validity}

The operator encoding applies in regimes where:
\begin{enumerate}[label=(\roman*)]
  \item a coherent spectral sector is isolated;
  \item the discarded sector has a positive damping gap;
  \item the admissible proper-time scale is finite;
  \item the retained sector supports coherent phase evolution;
  \item detector effects form a valid POVM or a clearly specified open measurement;
  \item a quantum instrument supplies the operational state update;
  \item branch damping defines a completely positive trace-preserving channel;
  \item survivor-basin capture is well-defined for the measurement context.
\end{enumerate}

Within this domain, wave-like interference and particle-like localization admit one consistent
operator encoding. Their derivation from one MTT source additionally requires intertwiners
from the source to the quantum state, dynamics, effects, instrument, and basin data.

Outside this domain, the theorem does not apply directly. The failure mode itself then becomes
physically meaningful: it indicates loss of coherent representability, gap closure, invalid
measurement modeling, or transition into a boundary regime.

\subsection{Summary}
\label{subsec:failure-summary}

The projection-duality account is not a universal slogan. It is a conditional structural result.
Its strength comes from its explicit hypotheses. If the hypotheses hold, then the wave and
particle descriptions are recovered as shadows of one finite coherent-sector excitation. If the
hypotheses fail, the theory identifies why the simple wave-particle description breaks down.

\begin{center}
\fbox{\parbox{0.86\linewidth}{\centering
The conditional encoding is valid where coherent projection, phase transport, an
operational instrument, and the declared basin intertwiner are valid.}}
\end{center}
\section{Conclusion}
\label{sec:conclusion}

Wave--particle duality is usually presented as a primitive feature of quantum mechanics:
microscopic systems propagate and interfere as waves, yet appear as localized particles when
measured. Modal Triplet Theory reframes this duality as a projection phenomenon. The
wave-like and particle-like descriptions are not two competing ontologies. They are two
downstream shadows of one finite coherent-sector excitation.

The central object is the admissible coherent operator
\begin{equation}
  \Badm
  =
  P\chi(A)e^{-\ta A}\chi(A)P.
\end{equation}
In its local representation,
\begin{equation}
  \Kadm(x,y)=\langle x|\Badm|y\rangle,
\end{equation}
this operator defines a finite source, response, or detector-amplitude kernel. When its width
is below the resolution scale of the probing apparatus, it appears pointlike. In the sharp
spectral/proper-time limit, it recovers the Dirac delta:
\begin{equation}
  \Kadm(x,y)\longrightarrow\delta(x-y).
\end{equation}
This is the particle shadow.

In its spectral representation,
\begin{equation}
  \Kadm(x,y)
  =
  \sum_{n\in\coh}
  \chi(\lambda_n)^2e^{-\ta\lambda_n}
  \phi_n(x)\phi_n^\ast(y),
\end{equation}
the same operator selects retained coherent modes. When those modes are transported by a
self-adjoint coherent phase generator \(H_{\rm coh}\), their relative phases produce
interference. This is the wave shadow.

Measurement connects these shadows. A measurement device is not a passive reader of a
context-free classical property. It is a device-specific disturbance and stabilization context:
\begin{equation}
  \Meas
  =
  (A_{\Meas},P_{\Meas},\chi_{\Meas},\tau_{\Meas},
  \{E_i^{(\Meas)}\},\{\mathcal I_i^{(\Meas)}\},\mathfrak B_{\Meas}).
\end{equation}
Different devices can therefore define different branch bases, finite detector effects,
instruments, damping matrices, and proposed survivor-basin partitions. The instrument
defines the operational record; the declared intertwiner identifies it with a basin record.

The interference side is controlled by off-diagonal branch coherence. A measurement-induced
damping map has the form
\begin{equation}
  \rho_{ab}\longmapsto D_{ab}^{(\Meas)}\rho_{ab},
\end{equation}
or
\begin{equation}
  \rho\longmapsto D^{(\Meas)}\circ\rho.
\end{equation}
For this to be physically admissible, \(D\) must define a valid Schur channel:
\begin{equation}
  D\succeq0,
  \qquad
  D_{aa}=1.
\end{equation}
When
\begin{equation}
  D_{ab}\approx1,
\end{equation}
branch coherence survives and the wave-like shadow is visible. When
\begin{equation}
  D_{ab}\approx0,
\end{equation}
branch coherence is suppressed and the downstream description becomes particle-like.

The double-slit experiment then becomes transparent. With no which-way record stabilized,
the two path branches remain coherent and the screen distribution contains the interference
term. With a which-way device, the device introduces branch separation and damps the
off-diagonal term. In a symmetric two-slit arrangement,
\begin{equation}
  V_{\rm MTT}
  =
  D_{12}V_0.
\end{equation}
Each individual screen event is localized by a finite detector effect, while the ensemble
pattern reflects the survival or loss of branch coherence.

This paper also separates two issues that are often conflated. Basin-measure probabilities
answer:
\begin{equation}
  \text{Which survivor basin is selected?}
\end{equation}
Branch damping answers:
\begin{equation}
  \text{How much interference remains between alternatives?}
\end{equation}
Thus Born frequencies and interference visibility are related but distinct layers of the MTT
measurement architecture.

The resulting conditional statement is:
\begin{equation}
  \boxed{
  \text{wave--particle duality}
  \quad\leadsto\quad
  \text{a projection-duality encoding}.
  }
\end{equation}
More explicitly,
\begin{equation}
  \boxed{
  \text{one finite coherent-sector excitation}
  \quad\Longrightarrow\quad
  \begin{cases}
  \text{local delta-like particle shadow},\\
  \text{spectral phase-coherent wave shadow},\\
  \text{finite measurement-selection shadow}.
  \end{cases}
  }
\end{equation}

The ordinary quantum descriptions are recovered in the appropriate limits. Standard wave
mechanics is recovered when coherent phases are retained. Pointlike detection is recovered
when finite detector effects are idealized by sharp projectors. Incoherent mixtures are
recovered when branch coherences are strongly damped. Thus the encoding is compatible with
the standard formalism. It is not yet a derivation of that formalism from MTT.

The remaining execution tasks are sector-specific. To turn the scaling relations into numerical
predictions, one must derive or constrain:
\begin{equation}
  A,\quad P,\quad \chi,\quad \lambda_\ast,\quad C_Q,\quad \epsadm,
\end{equation}
and, for each measurement context,
\begin{equation}
  A_{\Meas},\quad P_{\Meas},\quad \chi_{\Meas},\quad \tau_{\Meas},\quad
  \Lambda_{ab}^{(\Meas)},\quad \{E_i^{(\Meas)}\},\quad
  \{\mathcal I_i^{(\Meas)}\},\quad \mathfrak B_{\Meas}.
\end{equation}
Only then does the theory produce numerical predictions for finite-width deviations,
visibility loss, detector-resolution effects, or threshold knees.

The structural result is therefore conditional but useful: one filtered operator can support
local and spectral representations compatible with localized and interference experiments.
Whether both descend from one physical MTT excitation remains the selected-source theorem
to be proved.

\begin{thebibliography}{20}

\bibitem{ReedSimon}
M. Reed and B. Simon,
\emph{Methods of Modern Mathematical Physics I: Functional Analysis},
Academic Press.

\bibitem{Davies}
E. B. Davies,
\emph{Heat Kernels and Spectral Theory},
Cambridge University Press.

\bibitem{GlimmJaffe}
J. Glimm and A. Jaffe,
\emph{Quantum Physics: A Functional Integral Point of View},
Springer.

\bibitem{Zurek}
W. H. Zurek,
Decoherence, einselection, and the quantum origins of the classical,
\emph{Reviews of Modern Physics} 75, 715--775.

\bibitem{Schlosshauer}
M. Schlosshauer,
\emph{Decoherence and the Quantum-to-Classical Transition},
Springer.

\bibitem{Bohr}
N. Bohr,
The quantum postulate and the recent development of atomic theory,
\emph{Nature} 121, 580--590.

\bibitem{Bohm}
D. Bohm,
A suggested interpretation of the quantum theory in terms of hidden variables,
\emph{Physical Review} 85, 166--193.

\bibitem{FeynmanHibbs}
R. P. Feynman and A. R. Hibbs,
\emph{Quantum Mechanics and Path Integrals},
McGraw--Hill.

\bibitem{MTTDelta}
P. Nero,
\emph{Dirac Delta Functions as Singular Shadows of Admissible Projection},
MTT Delta--Projection sequence.

\bibitem{MTTKernel}
P. Nero,
\emph{Projected Heat Kernels from MTT Fixed-Point Data},
MTT Delta--Projection sequence.

\bibitem{MTTDamping}
P. Nero,
\emph{Admissibility-Time Bounds from Fixed-Point Damping},
MTT Delta--Projection sequence.


\bibitem{Schoenberg}
I. J. Schoenberg,
Metric spaces and positive definite functions,
\emph{Transactions of the American Mathematical Society} 44, 522--536.

\bibitem{Englert}
B.-G. Englert,
Fringe visibility and which-way information: An inequality,
\emph{Physical Review Letters} 77, 2154--2157.

\bibitem{HornJohnson}
R. A. Horn and C. R. Johnson,
\emph{Matrix Analysis},
Cambridge University Press.

\bibitem{Paulsen}
V. Paulsen,
\emph{Completely Bounded Maps and Operator Algebras},
Cambridge University Press.

\bibitem{NielsenChuang}
M. A. Nielsen and I. L. Chuang,
\emph{Quantum Computation and Quantum Information},
Cambridge University Press.

\bibitem{DaviesLewis}
E. B. Davies and J. T. Lewis,
``An operational approach to quantum probability,''
\emph{Communications in Mathematical Physics} \textbf{17} (1970), 239--260.
\href{https://doi.org/10.1007/BF01647093}{doi:10.1007/BF01647093}.

\bibitem{Ozawa}
M. Ozawa,
``Quantum measuring processes of continuous observables,''
\emph{Journal of Mathematical Physics} \textbf{25} (1984), 79--87.
\href{https://doi.org/10.1063/1.526000}{doi:10.1063/1.526000}.

\bibitem{Polchinski}
J. Polchinski,
\emph{String Theory, Volumes 1 and 2},
Cambridge University Press.

\bibitem{MTTBornBridge}
P. Nero,
\emph{Inflationary Measures and the Born Rule as a Single Shadow--Bridge Problem},
MTT Shadow--Bridge sequence.

\bibitem{MTTDisturbance}
P. Nero,
\emph{Measurement as Disturbance and Stabilization in Modal Triplet Theory},
MTT Measurement sequence.

\bibitem{MTTMeasurement}
P. Nero,
\emph{Measurement Effects as Finite Survivor-Basin Kernels},
MTT Delta--Projection sequence.

\end{thebibliography}

\end{document}
